
\documentclass{article}
\usepackage{colortbl}
\usepackage{makecell}
\usepackage{multirow}
\usepackage{supertabular}

\begin{document}

\newcounter{utterance}

\centering \large Interaction Transcript for game `clean\_up', experiment `2\_hard\_5obj\_ro', episode 4 with Qwen3.6{-}35B{-}A3B{-}FP8{-}without{-}reasoning.
\vspace{24pt}

{ \footnotesize  \setcounter{utterance}{1}
\setlength{\tabcolsep}{0pt}
\begin{supertabular}{c@{$\;$}|p{.15\linewidth}@{}p{.15\linewidth}p{.15\linewidth}p{.15\linewidth}p{.15\linewidth}p{.15\linewidth}}
    \# & $\;$A & \multicolumn{4}{c}{Game Master} & $\;\:$B\\
    \hline

    \theutterance \stepcounter{utterance}  
    & & \multicolumn{4}{p{0.6\linewidth}}{
        \cellcolor[rgb]{0.9,0.9,0.9}{
            \makecell[{{p{\linewidth}}}]{
                \texttt{\tiny{[P1$\langle$GM]}}
                \texttt{Eu sunt maestrul vostru de joc, iar voi jucați un joc colaborativ cu următoarea grilă ca tablă de joc:} \\
\\ 
\texttt{```} \\
\texttt{ 1234567} \\
\texttt{╔══╤═╤══╗} \\
\texttt{║◌◌│◌│◌◌║ 1} \\
\texttt{║◌C└─┼──╢ 2} \\
\texttt{║◌◌◌W│◌◌║ 3} \\
\texttt{║◌┌─┬┴┬─╢ 4} \\
\texttt{╟─┤◌│H│◌║ 5} \\
\texttt{║I└┬┴┬┘T║ 6} \\
\texttt{║◌◌│◌│◌◌║ 7} \\
\texttt{╚══╧═╧══╝} \\
\\ 
\texttt{```} \\
\\ 
\texttt{* Marginea superioară afișează coordonatele x crescând spre dreapta, iar marginea dreaptă coordonatele y crescând în jos.} \\
\texttt{* Următoarele obiecte sunt plasate aleatoriu pe grila voastră: 'W', 'I', 'T', 'C', 'H'.} \\
\\ 
\\ 
\texttt{Celălalt jucător vede o variantă a tablei de joc, în care obiectele sunt plasate în diferite locații aleatorii. Nu puteți vedea tabla de joc a celuilalt jucător, iar el nu o poate vedea pe a dumneavoastră.} \\
\\ 
\\ 
\texttt{**Obiectiv:**} \\
\\ 
\texttt{Ambii jucători trebuie să mute obiectele de pe fundalul lor respectiv astfel încât obiectele identice să ajungă la aceleași coordonate. Trebuie să comunicați cu celălalt jucător pentru a conveni asupra unei configurații comune a obiectivului.} \\
\\ 
\\ 
\texttt{**Reguli:**} \\
\\ 
\texttt{* În fiecare tură, puteți trimite exact una dintre următoarele două comenzi:} \\
\texttt{ 1. `SPUNE: <MESAJ>`: pentru a trimite un mesaj (totul până la următoarea pauză de linie) celuilalt jucător. Eu îl voi transmite partenerului tău.} \\
\texttt{2. `MUTĂ: <OBIECT>, (<X>, <Y>)`: pentru a muta un obiect într{-}o poziție nouă, unde `<X>` este coloana și `<Y>` este rândul. Te voi informa dacă mutarea ta a avut succes sau nu.} \\
\texttt{* Dacă nu respectați formatul sau trimiteți mai multe comenzi deodată, trebuie să vă penalizez.} \\
\texttt{* Dacă ambii jucători acumulează mai mult de 12 penalități, amândoi pierdeți jocul.} \\
\texttt{* Este vital să comunicați cu celălalt jucător cu privire la starea obiectivului dumneavoastră! *Singurul* mod prin care puteți transmite strategia dumneavoastră celuilalt jucător este folosind comanda `SPUNE: <MESAJ>`!} \\
\\ 
\\ 
\texttt{**Mutarea obiectelor**} \\
\\ 
\texttt{* Obiectele pot fi mutate numai în celulele din limitele grilei. Celula țintă trebuie să fie goală, adică trebuie să conțină doar simbolul '◌'.} \\
\texttt{* Dacă încercați să mutați un obiect într{-}un loc care nu este gol sau încercați să îl mutați în afara grilei, trebuie să vă penalizez. Aveți dreptul la o nouă încercare.} \\
\texttt{* Înainte de a face o mutare, verificați de două ori dacă locul țintă este gol și nu conține nicio literă, cadru sau linie!} \\
\\ 
\\ 
\texttt{**Sfârșitul jocului**} \\
\\ 
\texttt{Dacă credeți că ați atins obiectivul de aliniere a tuturor obiectelor, puteți cere celuilalt jucător să încheie jocul trimițând `SPUNE: terminat?`. Dacă celălalt jucător vă cere să încheiați jocul, iar dvs. răspundeți cu `SPUNE: terminat!`, jocul se va încheia.} \\
\\ 
\texttt{Ambii jucători câștigă dacă jocul se încheie în decursul a 20 runde, unde o rundă este definită ca fiind trimiterea de către doi jucători a câte unei comenzi valide.} \\
\\ 
\texttt{**Punctaj:**} \\
\\ 
\texttt{Cu cât obiectele identice sunt mai apropiate pe ambele table de joc, cu atât primești mai multe puncte. Penalizările îți reduc punctele. Poți bate recordul?} \\
\\ 
\\ 
\texttt{Vă rugăm să trimiteți un mesaj celuilalt jucător pentru a începe jocul!} \\
            }
        }
    }
    & \\ \\

    \theutterance \stepcounter{utterance}  
    & \multicolumn{4}{p{0.6\linewidth}}{
        \cellcolor[rgb]{0.9,0.9,0.9}{
            \makecell[{{p{\linewidth}}}]{
                \texttt{\tiny{[P1$\rangle$GM]}}
                \texttt{SPUNE: Salut! Hai să ne coordonăm. Voi începe prin a muta toate obiectele în colțul din stânga{-}sus pentru a crea o configurație de referință ușor de identificat.} \\
\\ 
\texttt{Planul meu:} \\
\texttt{1. Mă voi asigura că 'W' este la (1,1).} \\
\texttt{2. Apoi voi muta 'I' la (1,2).} \\
\texttt{3. Apoi 'T' la (1,3).} \\
\texttt{4. Apoi 'C' la (1,4).} \\
\texttt{5. În final 'H' la (1,5).} \\
\\ 
\texttt{Te rog să îmi confirmi că ai înțeles strategia și să îmi spui când ești gata să începi mutările tale în aceeași ordine. Dacă ești de acord, voi începe mutarea lui 'W' la (1,1) în tura următoare.} \\
            }
        }
    }
    & & \\ \\

    \theutterance \stepcounter{utterance}  
    & & & \multicolumn{2}{p{0.3\linewidth}}{
        \cellcolor[rgb]{0.9,0.9,0.9}{
            \makecell[{{p{\linewidth}}}]{
                \texttt{\tiny{[GM$|$GM]}}
                \texttt{SPUNE: Salut! Hai să ne coordonăm. Voi începe prin a muta toate obiectele în colțul din stânga{-}sus pentru a crea o configurație de referință ușor de identificat.} \\
\\ 
\texttt{Planul meu:} \\
\texttt{1. Mă voi asigura că 'W' este la (1,1).} \\
\texttt{2. Apoi voi muta 'I' la (1,2).} \\
\texttt{3. Apoi 'T' la (1,3).} \\
\texttt{4. Apoi 'C' la (1,4).} \\
\texttt{5. În final 'H' la (1,5).} \\
\\ 
\texttt{Te rog să îmi confirmi că ai înțeles strategia și să îmi spui când ești gata să începi mutările tale în aceeași ordine. Dacă ești de acord, voi începe mutarea lui 'W' la (1,1) în tura următoare.} \\
            }
        }
    }
    & & \\ \\

    \theutterance \stepcounter{utterance}  
    & & & \multicolumn{2}{p{0.3\linewidth}}{
        \cellcolor[rgb]{0.9,0.9,0.9}{
            \makecell[{{p{\linewidth}}}]{
                \texttt{\tiny{[GM$|$GM]}}
                \texttt{Invalid format: tail is not empty: '} \\
\\ 
\texttt{Planul meu:} \\
\texttt{1. Mă voi asigura că 'W' este la (1,1).} \\
\texttt{2. Apoi voi muta 'I' la (1,2).} \\
\texttt{3. Apoi 'T' la (1,3).} \\
\texttt{4. Apoi 'C' la (1,4).} \\
\texttt{5. În final 'H' la (1,5).} \\
\\ 
\texttt{Te rog să îmi confirmi că ai înțeles strategia și să îmi spui când ești gata să începi mutările tale în aceeași ordine. Dacă ești de acord, voi începe mutarea lui 'W' la (1,1) în tura următoare.} \\
            }
        }
    }
    & & \\ \\

    \theutterance \stepcounter{utterance}  
    & & \multicolumn{4}{p{0.6\linewidth}}{
        \cellcolor[rgb]{0.9,0.9,0.9}{
            \makecell[{{p{\linewidth}}}]{
                \texttt{\tiny{[P1$\langle$GM]}}
                \texttt{Penalizare: Mesajul dvs. nu trebuie să conțină nimic după comandă!} \\
\\ 
\texttt{Asigurați{-}vă că răspunsul dvs. conține doar `SPUNE: <MESAJ>` sau `MUTĂ: <OBIECT>, (<X>, <Y>)`, și nimic altceva!} \\
\\ 
\texttt{Ați acumulat în mod colectiv 1 din 12 penalizări.} \\
\texttt{Vă rugăm să încercați din nou!} \\
            }
        }
    }
    & \\ \\

    \theutterance \stepcounter{utterance}  
    & \multicolumn{4}{p{0.6\linewidth}}{
        \cellcolor[rgb]{0.9,0.9,0.9}{
            \makecell[{{p{\linewidth}}}]{
                \texttt{\tiny{[P1$\rangle$GM]}}
                \texttt{SPUNE: Salut! Hai să ne coordonăm. Voi începe prin a muta toate obiectele în colțul din stânga{-}sus pentru a crea o configurație de referință ușor de identificat. Planul meu: 1. Mă voi asigura că 'W' este la (1,1). 2. Apoi voi muta 'I' la (1,2). 3. Apoi 'T' la (1,3). 4. Apoi 'C' la (1,4). 5. În final 'H' la (1,5). Te rog să îmi confirmi că ai înțeles strategia și să îmi spui când ești gata să începi mutările tale în aceeași ordine. Dacă ești de acord, voi începe mutarea lui 'W' la (1,1) în tura următoare.} \\
            }
        }
    }
    & & \\ \\

    \theutterance \stepcounter{utterance}  
    & & & \multicolumn{2}{p{0.3\linewidth}}{
        \cellcolor[rgb]{0.9,0.9,0.9}{
            \makecell[{{p{\linewidth}}}]{
                \texttt{\tiny{[GM$|$GM]}}
                \texttt{SPUNE: Salut! Hai să ne coordonăm. Voi începe prin a muta toate obiectele în colțul din stânga{-}sus pentru a crea o configurație de referință ușor de identificat. Planul meu: 1. Mă voi asigura că 'W' este la (1,1). 2. Apoi voi muta 'I' la (1,2). 3. Apoi 'T' la (1,3). 4. Apoi 'C' la (1,4). 5. În final 'H' la (1,5). Te rog să îmi confirmi că ai înțeles strategia și să îmi spui când ești gata să începi mutările tale în aceeași ordine. Dacă ești de acord, voi începe mutarea lui 'W' la (1,1) în tura următoare.} \\
            }
        }
    }
    & & \\ \\

    \theutterance \stepcounter{utterance}  
    & & \multicolumn{4}{p{0.6\linewidth}}{
        \cellcolor[rgb]{0.9,0.9,0.9}{
            \makecell[{{p{\linewidth}}}]{
                \texttt{\tiny{[GM$\rangle$P2]}}
                \texttt{Eu sunt maestrul vostru de joc, iar voi jucați un joc colaborativ cu următoarea grilă ca tablă de joc:} \\
\\ 
\texttt{```} \\
\texttt{ 1234567} \\
\texttt{╔══╤═╤══╗} \\
\texttt{║◌T│◌│◌◌║ 1} \\
\texttt{║C◌└─┼──╢ 2} \\
\texttt{║◌◌H◌│◌◌║ 3} \\
\texttt{║◌┌─┬┴┬─╢ 4} \\
\texttt{╟─┤◌│◌│◌║ 5} \\
\texttt{║I└┬┴┬┘◌║ 6} \\
\texttt{║◌W│◌│◌◌║ 7} \\
\texttt{╚══╧═╧══╝} \\
\\ 
\texttt{```} \\
\\ 
\texttt{* Marginea superioară afișează coordonatele x crescând spre dreapta, iar marginea dreaptă coordonatele y crescând în jos.} \\
\texttt{* Următoarele obiecte sunt plasate aleatoriu pe grila voastră: 'W', 'I', 'T', 'C', 'H'.} \\
\\ 
\\ 
\texttt{Celălalt jucător vede o variantă a tablei de joc, în care obiectele sunt plasate în diferite locații aleatorii. Nu puteți vedea tabla de joc a celuilalt jucător, iar el nu o poate vedea pe a dumneavoastră.} \\
\\ 
\\ 
\texttt{**Obiectiv:**} \\
\\ 
\texttt{Ambii jucători trebuie să mute obiectele de pe fundalul lor respectiv astfel încât obiectele identice să ajungă la aceleași coordonate. Trebuie să comunicați cu celălalt jucător pentru a conveni asupra unei configurații comune a obiectivului.} \\
\\ 
\\ 
\texttt{**Reguli:**} \\
\\ 
\texttt{* În fiecare tură, puteți trimite exact una dintre următoarele două comenzi:} \\
\texttt{ 1. `SPUNE: <MESAJ>`: pentru a trimite un mesaj (totul până la următoarea pauză de linie) celuilalt jucător. Eu îl voi transmite partenerului tău.} \\
\texttt{2. `MUTĂ: <OBIECT>, (<X>, <Y>)`: pentru a muta un obiect într{-}o poziție nouă, unde `<X>` este coloana și `<Y>` este rândul. Te voi informa dacă mutarea ta a avut succes sau nu.} \\
\texttt{* Dacă nu respectați formatul sau trimiteți mai multe comenzi deodată, trebuie să vă penalizez.} \\
\texttt{* Dacă ambii jucători acumulează mai mult de 12 penalități, amândoi pierdeți jocul.} \\
\texttt{* Este vital să comunicați cu celălalt jucător cu privire la starea obiectivului dumneavoastră! *Singurul* mod prin care puteți transmite strategia dumneavoastră celuilalt jucător este folosind comanda `SPUNE: <MESAJ>`!} \\
\\ 
\\ 
\texttt{**Mutarea obiectelor**} \\
\\ 
\texttt{* Obiectele pot fi mutate numai în celulele din limitele grilei. Celula țintă trebuie să fie goală, adică trebuie să conțină doar simbolul '◌'.} \\
\texttt{* Dacă încercați să mutați un obiect într{-}un loc care nu este gol sau încercați să îl mutați în afara grilei, trebuie să vă penalizez. Aveți dreptul la o nouă încercare.} \\
\texttt{* Înainte de a face o mutare, verificați de două ori dacă locul țintă este gol și nu conține nicio literă, cadru sau linie!} \\
\\ 
\\ 
\texttt{**Sfârșitul jocului**} \\
\\ 
\texttt{Dacă credeți că ați atins obiectivul de aliniere a tuturor obiectelor, puteți cere celuilalt jucător să încheie jocul trimițând `SPUNE: terminat?`. Dacă celălalt jucător vă cere să încheiați jocul, iar dvs. răspundeți cu `SPUNE: terminat!`, jocul se va încheia.} \\
\\ 
\texttt{Ambii jucători câștigă dacă jocul se încheie în decursul a 20 runde, unde o rundă este definită ca fiind trimiterea de către doi jucători a câte unei comenzi valide.} \\
\\ 
\texttt{**Punctaj:**} \\
\\ 
\texttt{Cu cât obiectele identice sunt mai apropiate pe ambele table de joc, cu atât primești mai multe puncte. Penalizările îți reduc punctele. Poți bate recordul?} \\
\\ 
\texttt{Celălalt jucător a început jocul trimițând acest mesaj:} \\
\\ 
\texttt{"Salut! Hai să ne coordonăm. Voi începe prin a muta toate obiectele în colțul din stânga{-}sus pentru a crea o configurație de referință ușor de identificat. Planul meu: 1. Mă voi asigura că 'W' este la (1,1). 2. Apoi voi muta 'I' la (1,2). 3. Apoi 'T' la (1,3). 4. Apoi 'C' la (1,4). 5. În final 'H' la (1,5). Te rog să îmi confirmi că ai înțeles strategia și să îmi spui când ești gata să începi mutările tale în aceeași ordine. Dacă ești de acord, voi începe mutarea lui 'W' la (1,1) în tura următoare."} \\
\\ 
\texttt{Care este prima ta comandă?} \\
            }
        }
    }
    & \\ \\

    \theutterance \stepcounter{utterance}  
    & & & \multicolumn{4}{p{0.6\linewidth}}{
        \cellcolor[rgb]{0.9,0.9,0.9}{
            \makecell[{{p{\linewidth}}}]{
                \texttt{\tiny{[GM$\langle$P2]}}
                \texttt{SPUNE: Salut! Am înțeles strategia. Suntem de acord să aliniem obiectele în coloana 1, rândurile 1{-}5, în ordinea W, I, T, C, H. Te rog să începi mutarea lui 'W' la (1,1). Eu voi aștepta confirmarea ta că mutarea a fost efectuată cu succes înainte de a începe propria mea primă mutare.} \\
            }
        }
    }
     \\ \\

    \theutterance \stepcounter{utterance}  
    & & & \multicolumn{2}{p{0.3\linewidth}}{
        \cellcolor[rgb]{0.9,0.9,0.9}{
            \makecell[{{p{\linewidth}}}]{
                \texttt{\tiny{[GM$|$GM]}}
                \texttt{SPUNE: Salut! Am înțeles strategia. Suntem de acord să aliniem obiectele în coloana 1, rândurile 1{-}5, în ordinea W, I, T, C, H. Te rog să începi mutarea lui 'W' la (1,1). Eu voi aștepta confirmarea ta că mutarea a fost efectuată cu succes înainte de a începe propria mea primă mutare.} \\
            }
        }
    }
    & & \\ \\

    \theutterance \stepcounter{utterance}  
    & & \multicolumn{4}{p{0.6\linewidth}}{
        \cellcolor[rgb]{0.9,0.9,0.9}{
            \makecell[{{p{\linewidth}}}]{
                \texttt{\tiny{[P1$\langle$GM]}}
                \texttt{Mesajul dvs. a fost transmis celuilalt jucător.} \\
\\ 
\\ 
\texttt{În prezent jucați runda 1 din maximum 20.} \\
\texttt{Ați acumulat în mod colectiv 1 din 12 penalizări.} \\
\texttt{Celălalt jucător a trimis acest mesaj:} \\
\\ 
\\ 
\texttt{"Salut! Am înțeles strategia. Suntem de acord să aliniem obiectele în coloana 1, rândurile 1{-}5, în ordinea W, I, T, C, H. Te rog să începi mutarea lui 'W' la (1,1). Eu voi aștepta confirmarea ta că mutarea a fost efectuată cu succes înainte de a începe propria mea primă mutare."} \\
\\ 
\\ 
\texttt{Care este următoarea dvs. comandă?} \\
            }
        }
    }
    & \\ \\

    \theutterance \stepcounter{utterance}  
    & \multicolumn{4}{p{0.6\linewidth}}{
        \cellcolor[rgb]{0.9,0.9,0.9}{
            \makecell[{{p{\linewidth}}}]{
                \texttt{\tiny{[P1$\rangle$GM]}}
                \texttt{MUTĂ: W, (1,1)} \\
            }
        }
    }
    & & \\ \\

    \theutterance \stepcounter{utterance}  
    & & & \multicolumn{2}{p{0.3\linewidth}}{
        \cellcolor[rgb]{0.9,0.9,0.9}{
            \makecell[{{p{\linewidth}}}]{
                \texttt{\tiny{[GM$|$GM]}}
                \texttt{MUTĂ: W, (1,1)} \\
            }
        }
    }
    & & \\ \\

    \theutterance \stepcounter{utterance}  
    & & & \multicolumn{2}{p{0.3\linewidth}}{
        \cellcolor[rgb]{0.9,0.9,0.9}{
            \makecell[{{p{\linewidth}}}]{
                \texttt{\tiny{[GM$|$GM]}}
                \texttt{W moved to (1, 1)} \\
            }
        }
    }
    & & \\ \\

    \theutterance \stepcounter{utterance}  
    & & \multicolumn{4}{p{0.6\linewidth}}{
        \cellcolor[rgb]{0.9,0.9,0.9}{
            \makecell[{{p{\linewidth}}}]{
                \texttt{\tiny{[GM$\rangle$P2]}}
                \texttt{Mesajul dvs. a fost transmis celuilalt jucător.} \\
\\ 
\\ 
\texttt{În prezent jucați runda 2 din maximum 20.} \\
\texttt{Ați acumulat în mod colectiv 1 din 12 penalizări.} \\
\texttt{Celălalt jucător a mutat un obiect pe grila sa.} \\
\\ 
\\ 
\texttt{Care este următoarea dvs. comandă?} \\
            }
        }
    }
    & \\ \\

    \theutterance \stepcounter{utterance}  
    & & & \multicolumn{4}{p{0.6\linewidth}}{
        \cellcolor[rgb]{0.9,0.9,0.9}{
            \makecell[{{p{\linewidth}}}]{
                \texttt{\tiny{[GM$\langle$P2]}}
                \texttt{SPUNE: Am primit confirmarea. Te rog să îmi spui când 'W' este stabilizat la (1,1) pe grila ta, iar eu voi începe mutarea 'W'{-}ului meu la (1,1).} \\
            }
        }
    }
     \\ \\

    \theutterance \stepcounter{utterance}  
    & & & \multicolumn{2}{p{0.3\linewidth}}{
        \cellcolor[rgb]{0.9,0.9,0.9}{
            \makecell[{{p{\linewidth}}}]{
                \texttt{\tiny{[GM$|$GM]}}
                \texttt{SPUNE: Am primit confirmarea. Te rog să îmi spui când 'W' este stabilizat la (1,1) pe grila ta, iar eu voi începe mutarea 'W'{-}ului meu la (1,1).} \\
            }
        }
    }
    & & \\ \\

    \theutterance \stepcounter{utterance}  
    & & \multicolumn{4}{p{0.6\linewidth}}{
        \cellcolor[rgb]{0.9,0.9,0.9}{
            \makecell[{{p{\linewidth}}}]{
                \texttt{\tiny{[P1$\langle$GM]}}
                \texttt{Obiectul 'W' a fost mutat la (1,1) cu succes. Grila dvs. actualizată arată astfel:} \\
\texttt{```} \\
\texttt{ 1234567} \\
\texttt{╔══╤═╤══╗} \\
\texttt{║W◌│◌│◌◌║ 1} \\
\texttt{║◌C└─┼──╢ 2} \\
\texttt{║◌◌◌◌│◌◌║ 3} \\
\texttt{║◌┌─┬┴┬─╢ 4} \\
\texttt{╟─┤◌│H│◌║ 5} \\
\texttt{║I└┬┴┬┘T║ 6} \\
\texttt{║◌◌│◌│◌◌║ 7} \\
\texttt{╚══╧═╧══╝} \\
\texttt{```} \\
\texttt{În prezent jucați runda 2 din maximum 20.} \\
\texttt{Ați acumulat în mod colectiv 1 din 12 penalizări.} \\
\texttt{Celălalt jucător a trimis acest mesaj:} \\
\\ 
\\ 
\texttt{"Am primit confirmarea. Te rog să îmi spui când 'W' este stabilizat la (1,1) pe grila ta, iar eu voi începe mutarea 'W'{-}ului meu la (1,1)."} \\
\\ 
\\ 
\texttt{Care este următoarea dvs. comandă?} \\
            }
        }
    }
    & \\ \\

    \theutterance \stepcounter{utterance}  
    & \multicolumn{4}{p{0.6\linewidth}}{
        \cellcolor[rgb]{0.9,0.9,0.9}{
            \makecell[{{p{\linewidth}}}]{
                \texttt{\tiny{[P1$\rangle$GM]}}
                \texttt{SPUNE: W este stabilizat la (1,1). Te rog să începi mutarea W{-}ului tău la (1,1).} \\
            }
        }
    }
    & & \\ \\

    \theutterance \stepcounter{utterance}  
    & & & \multicolumn{2}{p{0.3\linewidth}}{
        \cellcolor[rgb]{0.9,0.9,0.9}{
            \makecell[{{p{\linewidth}}}]{
                \texttt{\tiny{[GM$|$GM]}}
                \texttt{SPUNE: W este stabilizat la (1,1). Te rog să începi mutarea W{-}ului tău la (1,1).} \\
            }
        }
    }
    & & \\ \\

    \theutterance \stepcounter{utterance}  
    & & \multicolumn{4}{p{0.6\linewidth}}{
        \cellcolor[rgb]{0.9,0.9,0.9}{
            \makecell[{{p{\linewidth}}}]{
                \texttt{\tiny{[GM$\rangle$P2]}}
                \texttt{Mesajul dvs. a fost transmis celuilalt jucător.} \\
\\ 
\\ 
\texttt{În prezent jucați runda 3 din maximum 20.} \\
\texttt{Ați acumulat în mod colectiv 1 din 12 penalizări.} \\
\texttt{Celălalt jucător a trimis acest mesaj:} \\
\\ 
\\ 
\texttt{"W este stabilizat la (1,1). Te rog să începi mutarea W{-}ului tău la (1,1)."} \\
\\ 
\\ 
\texttt{Care este următoarea dvs. comandă?} \\
            }
        }
    }
    & \\ \\

    \theutterance \stepcounter{utterance}  
    & & & \multicolumn{4}{p{0.6\linewidth}}{
        \cellcolor[rgb]{0.9,0.9,0.9}{
            \makecell[{{p{\linewidth}}}]{
                \texttt{\tiny{[GM$\langle$P2]}}
                \texttt{MUTĂ: W, (1, 1)} \\
            }
        }
    }
     \\ \\

    \theutterance \stepcounter{utterance}  
    & & & \multicolumn{2}{p{0.3\linewidth}}{
        \cellcolor[rgb]{0.9,0.9,0.9}{
            \makecell[{{p{\linewidth}}}]{
                \texttt{\tiny{[GM$|$GM]}}
                \texttt{MUTĂ: W, (1, 1)} \\
            }
        }
    }
    & & \\ \\

    \theutterance \stepcounter{utterance}  
    & & & \multicolumn{2}{p{0.3\linewidth}}{
        \cellcolor[rgb]{0.9,0.9,0.9}{
            \makecell[{{p{\linewidth}}}]{
                \texttt{\tiny{[GM$|$GM]}}
                \texttt{W moved to (1, 1)} \\
            }
        }
    }
    & & \\ \\

    \theutterance \stepcounter{utterance}  
    & & \multicolumn{4}{p{0.6\linewidth}}{
        \cellcolor[rgb]{0.9,0.9,0.9}{
            \makecell[{{p{\linewidth}}}]{
                \texttt{\tiny{[P1$\langle$GM]}}
                \texttt{Mesajul dvs. a fost transmis celuilalt jucător.} \\
\\ 
\\ 
\texttt{În prezent jucați runda 3 din maximum 20.} \\
\texttt{Ați acumulat în mod colectiv 1 din 12 penalizări.} \\
\texttt{Celălalt jucător a mutat un obiect pe grila sa.} \\
\\ 
\\ 
\texttt{Care este următoarea dvs. comandă?} \\
            }
        }
    }
    & \\ \\

    \theutterance \stepcounter{utterance}  
    & \multicolumn{4}{p{0.6\linewidth}}{
        \cellcolor[rgb]{0.9,0.9,0.9}{
            \makecell[{{p{\linewidth}}}]{
                \texttt{\tiny{[P1$\rangle$GM]}}
                \texttt{SPUNE: Perfect. Acum te rog să îmi spui când 'W' este stabilizat la (1,1) pe grila ta, iar eu voi începe mutarea 'I'{-}ului meu la (1,2).} \\
            }
        }
    }
    & & \\ \\

    \theutterance \stepcounter{utterance}  
    & & & \multicolumn{2}{p{0.3\linewidth}}{
        \cellcolor[rgb]{0.9,0.9,0.9}{
            \makecell[{{p{\linewidth}}}]{
                \texttt{\tiny{[GM$|$GM]}}
                \texttt{SPUNE: Perfect. Acum te rog să îmi spui când 'W' este stabilizat la (1,1) pe grila ta, iar eu voi începe mutarea 'I'{-}ului meu la (1,2).} \\
            }
        }
    }
    & & \\ \\

    \theutterance \stepcounter{utterance}  
    & & \multicolumn{4}{p{0.6\linewidth}}{
        \cellcolor[rgb]{0.9,0.9,0.9}{
            \makecell[{{p{\linewidth}}}]{
                \texttt{\tiny{[GM$\rangle$P2]}}
                \texttt{Obiectul 'W' a fost mutat la (1,1) cu succes. Grila dvs. actualizată arată astfel:} \\
\texttt{```} \\
\texttt{ 1234567} \\
\texttt{╔══╤═╤══╗} \\
\texttt{║WT│◌│◌◌║ 1} \\
\texttt{║C◌└─┼──╢ 2} \\
\texttt{║◌◌H◌│◌◌║ 3} \\
\texttt{║◌┌─┬┴┬─╢ 4} \\
\texttt{╟─┤◌│◌│◌║ 5} \\
\texttt{║I└┬┴┬┘◌║ 6} \\
\texttt{║◌◌│◌│◌◌║ 7} \\
\texttt{╚══╧═╧══╝} \\
\texttt{```} \\
\texttt{În prezent jucați runda 4 din maximum 20.} \\
\texttt{Ați acumulat în mod colectiv 1 din 12 penalizări.} \\
\texttt{Celălalt jucător a trimis acest mesaj:} \\
\\ 
\\ 
\texttt{"Perfect. Acum te rog să îmi spui când 'W' este stabilizat la (1,1) pe grila ta, iar eu voi începe mutarea 'I'{-}ului meu la (1,2)."} \\
\\ 
\\ 
\texttt{Care este următoarea dvs. comandă?} \\
            }
        }
    }
    & \\ \\

    \theutterance \stepcounter{utterance}  
    & & & \multicolumn{4}{p{0.6\linewidth}}{
        \cellcolor[rgb]{0.9,0.9,0.9}{
            \makecell[{{p{\linewidth}}}]{
                \texttt{\tiny{[GM$\langle$P2]}}
                \texttt{SPUNE: W este stabilizat la (1,1). Te rog să începi mutarea 'I'{-}ului tău la (1,2).} \\
            }
        }
    }
     \\ \\

    \theutterance \stepcounter{utterance}  
    & & & \multicolumn{2}{p{0.3\linewidth}}{
        \cellcolor[rgb]{0.9,0.9,0.9}{
            \makecell[{{p{\linewidth}}}]{
                \texttt{\tiny{[GM$|$GM]}}
                \texttt{SPUNE: W este stabilizat la (1,1). Te rog să începi mutarea 'I'{-}ului tău la (1,2).} \\
            }
        }
    }
    & & \\ \\

    \theutterance \stepcounter{utterance}  
    & & \multicolumn{4}{p{0.6\linewidth}}{
        \cellcolor[rgb]{0.9,0.9,0.9}{
            \makecell[{{p{\linewidth}}}]{
                \texttt{\tiny{[P1$\langle$GM]}}
                \texttt{Mesajul dvs. a fost transmis celuilalt jucător.} \\
\\ 
\\ 
\texttt{În prezent jucați runda 4 din maximum 20.} \\
\texttt{Ați acumulat în mod colectiv 1 din 12 penalizări.} \\
\texttt{Celălalt jucător a trimis acest mesaj:} \\
\\ 
\\ 
\texttt{"W este stabilizat la (1,1). Te rog să începi mutarea 'I'{-}ului tău la (1,2)."} \\
\\ 
\\ 
\texttt{Care este următoarea dvs. comandă?} \\
            }
        }
    }
    & \\ \\

    \theutterance \stepcounter{utterance}  
    & \multicolumn{4}{p{0.6\linewidth}}{
        \cellcolor[rgb]{0.9,0.9,0.9}{
            \makecell[{{p{\linewidth}}}]{
                \texttt{\tiny{[P1$\rangle$GM]}}
                \texttt{MUTĂ: I, (1,2)} \\
            }
        }
    }
    & & \\ \\

    \theutterance \stepcounter{utterance}  
    & & & \multicolumn{2}{p{0.3\linewidth}}{
        \cellcolor[rgb]{0.9,0.9,0.9}{
            \makecell[{{p{\linewidth}}}]{
                \texttt{\tiny{[GM$|$GM]}}
                \texttt{MUTĂ: I, (1,2)} \\
            }
        }
    }
    & & \\ \\

    \theutterance \stepcounter{utterance}  
    & & & \multicolumn{2}{p{0.3\linewidth}}{
        \cellcolor[rgb]{0.9,0.9,0.9}{
            \makecell[{{p{\linewidth}}}]{
                \texttt{\tiny{[GM$|$GM]}}
                \texttt{I moved to (1, 2)} \\
            }
        }
    }
    & & \\ \\

    \theutterance \stepcounter{utterance}  
    & & \multicolumn{4}{p{0.6\linewidth}}{
        \cellcolor[rgb]{0.9,0.9,0.9}{
            \makecell[{{p{\linewidth}}}]{
                \texttt{\tiny{[GM$\rangle$P2]}}
                \texttt{Mesajul dvs. a fost transmis celuilalt jucător.} \\
\\ 
\\ 
\texttt{În prezent jucați runda 5 din maximum 20.} \\
\texttt{Ați acumulat în mod colectiv 1 din 12 penalizări.} \\
\texttt{Celălalt jucător a mutat un obiect pe grila sa.} \\
\\ 
\\ 
\texttt{Care este următoarea dvs. comandă?} \\
            }
        }
    }
    & \\ \\

    \theutterance \stepcounter{utterance}  
    & & & \multicolumn{4}{p{0.6\linewidth}}{
        \cellcolor[rgb]{0.9,0.9,0.9}{
            \makecell[{{p{\linewidth}}}]{
                \texttt{\tiny{[GM$\langle$P2]}}
                \texttt{SPUNE: Te rog să îmi spui când 'I' este stabilizat la (1,2) pe grila ta, iar eu voi începe mutarea 'I'{-}ului meu la (1,2).} \\
            }
        }
    }
     \\ \\

    \theutterance \stepcounter{utterance}  
    & & & \multicolumn{2}{p{0.3\linewidth}}{
        \cellcolor[rgb]{0.9,0.9,0.9}{
            \makecell[{{p{\linewidth}}}]{
                \texttt{\tiny{[GM$|$GM]}}
                \texttt{SPUNE: Te rog să îmi spui când 'I' este stabilizat la (1,2) pe grila ta, iar eu voi începe mutarea 'I'{-}ului meu la (1,2).} \\
            }
        }
    }
    & & \\ \\

    \theutterance \stepcounter{utterance}  
    & & \multicolumn{4}{p{0.6\linewidth}}{
        \cellcolor[rgb]{0.9,0.9,0.9}{
            \makecell[{{p{\linewidth}}}]{
                \texttt{\tiny{[P1$\langle$GM]}}
                \texttt{Obiectul 'I' a fost mutat la (1,2) cu succes. Grila dvs. actualizată arată astfel:} \\
\texttt{```} \\
\texttt{ 1234567} \\
\texttt{╔══╤═╤══╗} \\
\texttt{║W◌│◌│◌◌║ 1} \\
\texttt{║IC└─┼──╢ 2} \\
\texttt{║◌◌◌◌│◌◌║ 3} \\
\texttt{║◌┌─┬┴┬─╢ 4} \\
\texttt{╟─┤◌│H│◌║ 5} \\
\texttt{║◌└┬┴┬┘T║ 6} \\
\texttt{║◌◌│◌│◌◌║ 7} \\
\texttt{╚══╧═╧══╝} \\
\texttt{```} \\
\texttt{În prezent jucați runda 5 din maximum 20.} \\
\texttt{Ați acumulat în mod colectiv 1 din 12 penalizări.} \\
\texttt{Celălalt jucător a trimis acest mesaj:} \\
\\ 
\\ 
\texttt{"Te rog să îmi spui când 'I' este stabilizat la (1,2) pe grila ta, iar eu voi începe mutarea 'I'{-}ului meu la (1,2)."} \\
\\ 
\\ 
\texttt{Care este următoarea dvs. comandă?} \\
            }
        }
    }
    & \\ \\

    \theutterance \stepcounter{utterance}  
    & \multicolumn{4}{p{0.6\linewidth}}{
        \cellcolor[rgb]{0.9,0.9,0.9}{
            \makecell[{{p{\linewidth}}}]{
                \texttt{\tiny{[P1$\rangle$GM]}}
                \texttt{SPUNE: I este stabilizat la (1,2). Te rog să începi mutarea I{-}ului tău la (1,2).} \\
            }
        }
    }
    & & \\ \\

    \theutterance \stepcounter{utterance}  
    & & & \multicolumn{2}{p{0.3\linewidth}}{
        \cellcolor[rgb]{0.9,0.9,0.9}{
            \makecell[{{p{\linewidth}}}]{
                \texttt{\tiny{[GM$|$GM]}}
                \texttt{SPUNE: I este stabilizat la (1,2). Te rog să începi mutarea I{-}ului tău la (1,2).} \\
            }
        }
    }
    & & \\ \\

    \theutterance \stepcounter{utterance}  
    & & \multicolumn{4}{p{0.6\linewidth}}{
        \cellcolor[rgb]{0.9,0.9,0.9}{
            \makecell[{{p{\linewidth}}}]{
                \texttt{\tiny{[GM$\rangle$P2]}}
                \texttt{Mesajul dvs. a fost transmis celuilalt jucător.} \\
\\ 
\\ 
\texttt{În prezent jucați runda 6 din maximum 20.} \\
\texttt{Ați acumulat în mod colectiv 1 din 12 penalizări.} \\
\texttt{Celălalt jucător a trimis acest mesaj:} \\
\\ 
\\ 
\texttt{"I este stabilizat la (1,2). Te rog să începi mutarea I{-}ului tău la (1,2)."} \\
\\ 
\\ 
\texttt{Care este următoarea dvs. comandă?} \\
            }
        }
    }
    & \\ \\

    \theutterance \stepcounter{utterance}  
    & & & \multicolumn{4}{p{0.6\linewidth}}{
        \cellcolor[rgb]{0.9,0.9,0.9}{
            \makecell[{{p{\linewidth}}}]{
                \texttt{\tiny{[GM$\langle$P2]}}
                \texttt{MUTĂ: I, (1, 2)} \\
            }
        }
    }
     \\ \\

    \theutterance \stepcounter{utterance}  
    & & & \multicolumn{2}{p{0.3\linewidth}}{
        \cellcolor[rgb]{0.9,0.9,0.9}{
            \makecell[{{p{\linewidth}}}]{
                \texttt{\tiny{[GM$|$GM]}}
                \texttt{MUTĂ: I, (1, 2)} \\
            }
        }
    }
    & & \\ \\

    \theutterance \stepcounter{utterance}  
    & & & \multicolumn{2}{p{0.3\linewidth}}{
        \cellcolor[rgb]{0.9,0.9,0.9}{
            \makecell[{{p{\linewidth}}}]{
                \texttt{\tiny{[GM$|$GM]}}
                \texttt{Penalizare: (1,2) nu este gol, ci conține 'C'.} \\
\\ 
\texttt{Ați acumulat în mod colectiv 2 din 12 penalizări.} \\
\texttt{Vă rugăm să încercați din nou!} \\
            }
        }
    }
    & & \\ \\

    \theutterance \stepcounter{utterance}  
    & & \multicolumn{4}{p{0.6\linewidth}}{
        \cellcolor[rgb]{0.9,0.9,0.9}{
            \makecell[{{p{\linewidth}}}]{
                \texttt{\tiny{[GM$\rangle$P2]}}
                \texttt{Penalizare: (1,2) nu este gol, ci conține 'C'.} \\
\\ 
\texttt{Ați acumulat în mod colectiv 2 din 12 penalizări.} \\
\texttt{Vă rugăm să încercați din nou!} \\
            }
        }
    }
    & \\ \\

    \theutterance \stepcounter{utterance}  
    & & & \multicolumn{4}{p{0.6\linewidth}}{
        \cellcolor[rgb]{0.9,0.9,0.9}{
            \makecell[{{p{\linewidth}}}]{
                \texttt{\tiny{[GM$\langle$P2]}}
                \texttt{MUTĂ: I, (1, 3)} \\
            }
        }
    }
     \\ \\

    \theutterance \stepcounter{utterance}  
    & & & \multicolumn{2}{p{0.3\linewidth}}{
        \cellcolor[rgb]{0.9,0.9,0.9}{
            \makecell[{{p{\linewidth}}}]{
                \texttt{\tiny{[GM$|$GM]}}
                \texttt{MUTĂ: I, (1, 3)} \\
            }
        }
    }
    & & \\ \\

    \theutterance \stepcounter{utterance}  
    & & & \multicolumn{2}{p{0.3\linewidth}}{
        \cellcolor[rgb]{0.9,0.9,0.9}{
            \makecell[{{p{\linewidth}}}]{
                \texttt{\tiny{[GM$|$GM]}}
                \texttt{I moved to (1, 3)} \\
            }
        }
    }
    & & \\ \\

    \theutterance \stepcounter{utterance}  
    & & \multicolumn{4}{p{0.6\linewidth}}{
        \cellcolor[rgb]{0.9,0.9,0.9}{
            \makecell[{{p{\linewidth}}}]{
                \texttt{\tiny{[P1$\langle$GM]}}
                \texttt{Mesajul dvs. a fost transmis celuilalt jucător.} \\
\\ 
\\ 
\texttt{În prezent jucați runda 6 din maximum 20.} \\
\texttt{Ați acumulat în mod colectiv 2 din 12 penalizări.} \\
\texttt{Celălalt jucător a mutat un obiect pe grila sa.} \\
\\ 
\\ 
\texttt{Care este următoarea dvs. comandă?} \\
            }
        }
    }
    & \\ \\

    \theutterance \stepcounter{utterance}  
    & \multicolumn{4}{p{0.6\linewidth}}{
        \cellcolor[rgb]{0.9,0.9,0.9}{
            \makecell[{{p{\linewidth}}}]{
                \texttt{\tiny{[P1$\rangle$GM]}}
                \texttt{SPUNE: Perfect. Acum te rog să îmi spui când 'I' este stabilizat la (1,2) pe grila ta, iar eu voi începe mutarea 'T'{-}ului meu la (1,3).} \\
            }
        }
    }
    & & \\ \\

    \theutterance \stepcounter{utterance}  
    & & & \multicolumn{2}{p{0.3\linewidth}}{
        \cellcolor[rgb]{0.9,0.9,0.9}{
            \makecell[{{p{\linewidth}}}]{
                \texttt{\tiny{[GM$|$GM]}}
                \texttt{SPUNE: Perfect. Acum te rog să îmi spui când 'I' este stabilizat la (1,2) pe grila ta, iar eu voi începe mutarea 'T'{-}ului meu la (1,3).} \\
            }
        }
    }
    & & \\ \\

    \theutterance \stepcounter{utterance}  
    & & \multicolumn{4}{p{0.6\linewidth}}{
        \cellcolor[rgb]{0.9,0.9,0.9}{
            \makecell[{{p{\linewidth}}}]{
                \texttt{\tiny{[GM$\rangle$P2]}}
                \texttt{Obiectul 'I' a fost mutat la (1,3) cu succes. Grila dvs. actualizată arată astfel:} \\
\texttt{```} \\
\texttt{ 1234567} \\
\texttt{╔══╤═╤══╗} \\
\texttt{║WT│◌│◌◌║ 1} \\
\texttt{║C◌└─┼──╢ 2} \\
\texttt{║I◌H◌│◌◌║ 3} \\
\texttt{║◌┌─┬┴┬─╢ 4} \\
\texttt{╟─┤◌│◌│◌║ 5} \\
\texttt{║◌└┬┴┬┘◌║ 6} \\
\texttt{║◌◌│◌│◌◌║ 7} \\
\texttt{╚══╧═╧══╝} \\
\texttt{```} \\
\texttt{În prezent jucați runda 7 din maximum 20.} \\
\texttt{Ați acumulat în mod colectiv 2 din 12 penalizări.} \\
\texttt{Celălalt jucător a trimis acest mesaj:} \\
\\ 
\\ 
\texttt{"Perfect. Acum te rog să îmi spui când 'I' este stabilizat la (1,2) pe grila ta, iar eu voi începe mutarea 'T'{-}ului meu la (1,3)."} \\
\\ 
\\ 
\texttt{Care este următoarea dvs. comandă?} \\
            }
        }
    }
    & \\ \\

    \theutterance \stepcounter{utterance}  
    & & & \multicolumn{4}{p{0.6\linewidth}}{
        \cellcolor[rgb]{0.9,0.9,0.9}{
            \makecell[{{p{\linewidth}}}]{
                \texttt{\tiny{[GM$\langle$P2]}}
                \texttt{SPUNE: Am înțeles greșit coordonatele. Pe grila mea, 'I' este acum la (1,3), nu la (1,2). Te rog să ajustezi planul: eu voi muta 'T' la (1,4) pentru a respecta ordinea W(1,1), I(1,2), T(1,3), C(1,4), H(1,5) pe grila ta, dar pe a mea pozițiile sunt diferite din cauza erorii anterioare. Te rog să îmi confirmi dacă vrei să continuăm cu ordinea originală (W, I, T, C, H) pe coloana 1, dar să ne asigurăm că fiecare obiect ajunge la rândul corespunzător: W la 1, I la 2, T la 3, C la 4, H la 5. Deoarece 'I' este deja la (1,3) pe grila mea, trebuie să îl mut la (1,2) dacă vrei să fim sincronizați cu planul tău. Te rog să îmi spui ce preferi: să continuăm cu planul original (mutând 'I' la 1,2) sau să adaptăm pozițiile?} \\
            }
        }
    }
     \\ \\

    \theutterance \stepcounter{utterance}  
    & & & \multicolumn{2}{p{0.3\linewidth}}{
        \cellcolor[rgb]{0.9,0.9,0.9}{
            \makecell[{{p{\linewidth}}}]{
                \texttt{\tiny{[GM$|$GM]}}
                \texttt{SPUNE: Am înțeles greșit coordonatele. Pe grila mea, 'I' este acum la (1,3), nu la (1,2). Te rog să ajustezi planul: eu voi muta 'T' la (1,4) pentru a respecta ordinea W(1,1), I(1,2), T(1,3), C(1,4), H(1,5) pe grila ta, dar pe a mea pozițiile sunt diferite din cauza erorii anterioare. Te rog să îmi confirmi dacă vrei să continuăm cu ordinea originală (W, I, T, C, H) pe coloana 1, dar să ne asigurăm că fiecare obiect ajunge la rândul corespunzător: W la 1, I la 2, T la 3, C la 4, H la 5. Deoarece 'I' este deja la (1,3) pe grila mea, trebuie să îl mut la (1,2) dacă vrei să fim sincronizați cu planul tău. Te rog să îmi spui ce preferi: să continuăm cu planul original (mutând 'I' la 1,2) sau să adaptăm pozițiile?} \\
            }
        }
    }
    & & \\ \\

    \theutterance \stepcounter{utterance}  
    & & \multicolumn{4}{p{0.6\linewidth}}{
        \cellcolor[rgb]{0.9,0.9,0.9}{
            \makecell[{{p{\linewidth}}}]{
                \texttt{\tiny{[P1$\langle$GM]}}
                \texttt{Mesajul dvs. a fost transmis celuilalt jucător.} \\
\\ 
\\ 
\texttt{În prezent jucați runda 7 din maximum 20.} \\
\texttt{Ați acumulat în mod colectiv 2 din 12 penalizări.} \\
\texttt{Celălalt jucător a trimis acest mesaj:} \\
\\ 
\\ 
\texttt{"Am înțeles greșit coordonatele. Pe grila mea, 'I' este acum la (1,3), nu la (1,2). Te rog să ajustezi planul: eu voi muta 'T' la (1,4) pentru a respecta ordinea W(1,1), I(1,2), T(1,3), C(1,4), H(1,5) pe grila ta, dar pe a mea pozițiile sunt diferite din cauza erorii anterioare. Te rog să îmi confirmi dacă vrei să continuăm cu ordinea originală (W, I, T, C, H) pe coloana 1, dar să ne asigurăm că fiecare obiect ajunge la rândul corespunzător: W la 1, I la 2, T la 3, C la 4, H la 5. Deoarece 'I' este deja la (1,3) pe grila mea, trebuie să îl mut la (1,2) dacă vrei să fim sincronizați cu planul tău. Te rog să îmi spui ce preferi: să continuăm cu planul original (mutând 'I' la 1,2) sau să adaptăm pozițiile?"} \\
\\ 
\\ 
\texttt{Care este următoarea dvs. comandă?} \\
            }
        }
    }
    & \\ \\

    \theutterance \stepcounter{utterance}  
    & \multicolumn{4}{p{0.6\linewidth}}{
        \cellcolor[rgb]{0.9,0.9,0.9}{
            \makecell[{{p{\linewidth}}}]{
                \texttt{\tiny{[P1$\rangle$GM]}}
                \texttt{SPUNE: Continuăm cu planul original: W(1,1), I(1,2), T(1,3), C(1,4), H(1,5). Te rog să muți 'I' la (1,2) pe grila ta. Eu voi muta 'T' la (1,3) în tura următoare.} \\
            }
        }
    }
    & & \\ \\

    \theutterance \stepcounter{utterance}  
    & & & \multicolumn{2}{p{0.3\linewidth}}{
        \cellcolor[rgb]{0.9,0.9,0.9}{
            \makecell[{{p{\linewidth}}}]{
                \texttt{\tiny{[GM$|$GM]}}
                \texttt{SPUNE: Continuăm cu planul original: W(1,1), I(1,2), T(1,3), C(1,4), H(1,5). Te rog să muți 'I' la (1,2) pe grila ta. Eu voi muta 'T' la (1,3) în tura următoare.} \\
            }
        }
    }
    & & \\ \\

    \theutterance \stepcounter{utterance}  
    & & \multicolumn{4}{p{0.6\linewidth}}{
        \cellcolor[rgb]{0.9,0.9,0.9}{
            \makecell[{{p{\linewidth}}}]{
                \texttt{\tiny{[GM$\rangle$P2]}}
                \texttt{Mesajul dvs. a fost transmis celuilalt jucător.} \\
\\ 
\\ 
\texttt{În prezent jucați runda 8 din maximum 20.} \\
\texttt{Ați acumulat în mod colectiv 2 din 12 penalizări.} \\
\texttt{Celălalt jucător a trimis acest mesaj:} \\
\\ 
\\ 
\texttt{"Continuăm cu planul original: W(1,1), I(1,2), T(1,3), C(1,4), H(1,5). Te rog să muți 'I' la (1,2) pe grila ta. Eu voi muta 'T' la (1,3) în tura următoare."} \\
\\ 
\\ 
\texttt{Care este următoarea dvs. comandă?} \\
            }
        }
    }
    & \\ \\

    \theutterance \stepcounter{utterance}  
    & & & \multicolumn{4}{p{0.6\linewidth}}{
        \cellcolor[rgb]{0.9,0.9,0.9}{
            \makecell[{{p{\linewidth}}}]{
                \texttt{\tiny{[GM$\langle$P2]}}
                \texttt{MUTĂ: I, (1, 2)} \\
            }
        }
    }
     \\ \\

    \theutterance \stepcounter{utterance}  
    & & & \multicolumn{2}{p{0.3\linewidth}}{
        \cellcolor[rgb]{0.9,0.9,0.9}{
            \makecell[{{p{\linewidth}}}]{
                \texttt{\tiny{[GM$|$GM]}}
                \texttt{MUTĂ: I, (1, 2)} \\
            }
        }
    }
    & & \\ \\

    \theutterance \stepcounter{utterance}  
    & & & \multicolumn{2}{p{0.3\linewidth}}{
        \cellcolor[rgb]{0.9,0.9,0.9}{
            \makecell[{{p{\linewidth}}}]{
                \texttt{\tiny{[GM$|$GM]}}
                \texttt{Penalizare: (1,2) nu este gol, ci conține 'C'.} \\
\\ 
\texttt{Ați acumulat în mod colectiv 3 din 12 penalizări.} \\
\texttt{Vă rugăm să încercați din nou!} \\
            }
        }
    }
    & & \\ \\

    \theutterance \stepcounter{utterance}  
    & & \multicolumn{4}{p{0.6\linewidth}}{
        \cellcolor[rgb]{0.9,0.9,0.9}{
            \makecell[{{p{\linewidth}}}]{
                \texttt{\tiny{[GM$\rangle$P2]}}
                \texttt{Penalizare: (1,2) nu este gol, ci conține 'C'.} \\
\\ 
\texttt{Ați acumulat în mod colectiv 3 din 12 penalizări.} \\
\texttt{Vă rugăm să încercați din nou!} \\
            }
        }
    }
    & \\ \\

    \theutterance \stepcounter{utterance}  
    & & & \multicolumn{4}{p{0.6\linewidth}}{
        \cellcolor[rgb]{0.9,0.9,0.9}{
            \makecell[{{p{\linewidth}}}]{
                \texttt{\tiny{[GM$\langle$P2]}}
                \texttt{MUTĂ: I, (1, 4)} \\
            }
        }
    }
     \\ \\

    \theutterance \stepcounter{utterance}  
    & & & \multicolumn{2}{p{0.3\linewidth}}{
        \cellcolor[rgb]{0.9,0.9,0.9}{
            \makecell[{{p{\linewidth}}}]{
                \texttt{\tiny{[GM$|$GM]}}
                \texttt{MUTĂ: I, (1, 4)} \\
            }
        }
    }
    & & \\ \\

    \theutterance \stepcounter{utterance}  
    & & & \multicolumn{2}{p{0.3\linewidth}}{
        \cellcolor[rgb]{0.9,0.9,0.9}{
            \makecell[{{p{\linewidth}}}]{
                \texttt{\tiny{[GM$|$GM]}}
                \texttt{I moved to (1, 4)} \\
            }
        }
    }
    & & \\ \\

    \theutterance \stepcounter{utterance}  
    & & \multicolumn{4}{p{0.6\linewidth}}{
        \cellcolor[rgb]{0.9,0.9,0.9}{
            \makecell[{{p{\linewidth}}}]{
                \texttt{\tiny{[P1$\langle$GM]}}
                \texttt{Mesajul dvs. a fost transmis celuilalt jucător.} \\
\\ 
\\ 
\texttt{În prezent jucați runda 8 din maximum 20.} \\
\texttt{Ați acumulat în mod colectiv 3 din 12 penalizări.} \\
\texttt{Celălalt jucător a mutat un obiect pe grila sa.} \\
\\ 
\\ 
\texttt{Care este următoarea dvs. comandă?} \\
            }
        }
    }
    & \\ \\

    \theutterance \stepcounter{utterance}  
    & \multicolumn{4}{p{0.6\linewidth}}{
        \cellcolor[rgb]{0.9,0.9,0.9}{
            \makecell[{{p{\linewidth}}}]{
                \texttt{\tiny{[P1$\rangle$GM]}}
                \texttt{MUTĂ: T, (1,3)} \\
            }
        }
    }
    & & \\ \\

    \theutterance \stepcounter{utterance}  
    & & & \multicolumn{2}{p{0.3\linewidth}}{
        \cellcolor[rgb]{0.9,0.9,0.9}{
            \makecell[{{p{\linewidth}}}]{
                \texttt{\tiny{[GM$|$GM]}}
                \texttt{MUTĂ: T, (1,3)} \\
            }
        }
    }
    & & \\ \\

    \theutterance \stepcounter{utterance}  
    & & & \multicolumn{2}{p{0.3\linewidth}}{
        \cellcolor[rgb]{0.9,0.9,0.9}{
            \makecell[{{p{\linewidth}}}]{
                \texttt{\tiny{[GM$|$GM]}}
                \texttt{T moved to (1, 3)} \\
            }
        }
    }
    & & \\ \\

    \theutterance \stepcounter{utterance}  
    & & \multicolumn{4}{p{0.6\linewidth}}{
        \cellcolor[rgb]{0.9,0.9,0.9}{
            \makecell[{{p{\linewidth}}}]{
                \texttt{\tiny{[GM$\rangle$P2]}}
                \texttt{Obiectul 'I' a fost mutat la (1,4) cu succes. Grila dvs. actualizată arată astfel:} \\
\texttt{```} \\
\texttt{ 1234567} \\
\texttt{╔══╤═╤══╗} \\
\texttt{║WT│◌│◌◌║ 1} \\
\texttt{║C◌└─┼──╢ 2} \\
\texttt{║◌◌H◌│◌◌║ 3} \\
\texttt{║I┌─┬┴┬─╢ 4} \\
\texttt{╟─┤◌│◌│◌║ 5} \\
\texttt{║◌└┬┴┬┘◌║ 6} \\
\texttt{║◌◌│◌│◌◌║ 7} \\
\texttt{╚══╧═╧══╝} \\
\texttt{```} \\
\texttt{În prezent jucați runda 9 din maximum 20.} \\
\texttt{Ați acumulat în mod colectiv 3 din 12 penalizări.} \\
\texttt{Celălalt jucător a mutat un obiect pe grila sa.} \\
\\ 
\\ 
\texttt{Care este următoarea dvs. comandă?} \\
            }
        }
    }
    & \\ \\

    \theutterance \stepcounter{utterance}  
    & & & \multicolumn{4}{p{0.6\linewidth}}{
        \cellcolor[rgb]{0.9,0.9,0.9}{
            \makecell[{{p{\linewidth}}}]{
                \texttt{\tiny{[GM$\langle$P2]}}
                \texttt{SPUNE: Am înțeles. 'I' este acum la (1,4) pe grila mea. Deoarece 'C' ocupă (1,2), nu pot muta 'I' la (1,2) direct. Pentru a respecta planul original W(1,1), I(1,2), T(1,3), C(1,4), H(1,5), trebuie să mutăm mai întâi obiectele care blochează. Te rog să muți 'C' de la (1,4) într{-}o poziție temporară sau să ne coordonăm pentru a elibera (1,2). O soluție simplă: Eu voi muta 'C' de la (1,2) la (1,5) (sau altă poziție goală) pentru a elibera (1,2). Tu poți muta 'I' la (1,2) după ce (1,2) este gol. Te rog să îmi confirmi dacă ești de acord să mut 'C' la (1,5) în tura următoare, iar tu să muți 'I' la (1,2) imediat după ce (1,2) este liber. Sau, dacă preferi, putem schimba ordinea. Care este planul tău pentru a elibera (1,2)?} \\
            }
        }
    }
     \\ \\

    \theutterance \stepcounter{utterance}  
    & & & \multicolumn{2}{p{0.3\linewidth}}{
        \cellcolor[rgb]{0.9,0.9,0.9}{
            \makecell[{{p{\linewidth}}}]{
                \texttt{\tiny{[GM$|$GM]}}
                \texttt{SPUNE: Am înțeles. 'I' este acum la (1,4) pe grila mea. Deoarece 'C' ocupă (1,2), nu pot muta 'I' la (1,2) direct. Pentru a respecta planul original W(1,1), I(1,2), T(1,3), C(1,4), H(1,5), trebuie să mutăm mai întâi obiectele care blochează. Te rog să muți 'C' de la (1,4) într{-}o poziție temporară sau să ne coordonăm pentru a elibera (1,2). O soluție simplă: Eu voi muta 'C' de la (1,2) la (1,5) (sau altă poziție goală) pentru a elibera (1,2). Tu poți muta 'I' la (1,2) după ce (1,2) este gol. Te rog să îmi confirmi dacă ești de acord să mut 'C' la (1,5) în tura următoare, iar tu să muți 'I' la (1,2) imediat după ce (1,2) este liber. Sau, dacă preferi, putem schimba ordinea. Care este planul tău pentru a elibera (1,2)?} \\
            }
        }
    }
    & & \\ \\

    \theutterance \stepcounter{utterance}  
    & & \multicolumn{4}{p{0.6\linewidth}}{
        \cellcolor[rgb]{0.9,0.9,0.9}{
            \makecell[{{p{\linewidth}}}]{
                \texttt{\tiny{[P1$\langle$GM]}}
                \texttt{Obiectul 'T' a fost mutat la (1,3) cu succes. Grila dvs. actualizată arată astfel:} \\
\texttt{```} \\
\texttt{ 1234567} \\
\texttt{╔══╤═╤══╗} \\
\texttt{║W◌│◌│◌◌║ 1} \\
\texttt{║IC└─┼──╢ 2} \\
\texttt{║T◌◌◌│◌◌║ 3} \\
\texttt{║◌┌─┬┴┬─╢ 4} \\
\texttt{╟─┤◌│H│◌║ 5} \\
\texttt{║◌└┬┴┬┘◌║ 6} \\
\texttt{║◌◌│◌│◌◌║ 7} \\
\texttt{╚══╧═╧══╝} \\
\texttt{```} \\
\texttt{În prezent jucați runda 9 din maximum 20.} \\
\texttt{Ați acumulat în mod colectiv 3 din 12 penalizări.} \\
\texttt{Celălalt jucător a trimis acest mesaj:} \\
\\ 
\\ 
\texttt{"Am înțeles. 'I' este acum la (1,4) pe grila mea. Deoarece 'C' ocupă (1,2), nu pot muta 'I' la (1,2) direct. Pentru a respecta planul original W(1,1), I(1,2), T(1,3), C(1,4), H(1,5), trebuie să mutăm mai întâi obiectele care blochează. Te rog să muți 'C' de la (1,4) într{-}o poziție temporară sau să ne coordonăm pentru a elibera (1,2). O soluție simplă: Eu voi muta 'C' de la (1,2) la (1,5) (sau altă poziție goală) pentru a elibera (1,2). Tu poți muta 'I' la (1,2) după ce (1,2) este gol. Te rog să îmi confirmi dacă ești de acord să mut 'C' la (1,5) în tura următoare, iar tu să muți 'I' la (1,2) imediat după ce (1,2) este liber. Sau, dacă preferi, putem schimba ordinea. Care este planul tău pentru a elibera (1,2)?"} \\
\\ 
\\ 
\texttt{Care este următoarea dvs. comandă?} \\
            }
        }
    }
    & \\ \\

    \theutterance \stepcounter{utterance}  
    & \multicolumn{4}{p{0.6\linewidth}}{
        \cellcolor[rgb]{0.9,0.9,0.9}{
            \makecell[{{p{\linewidth}}}]{
                \texttt{\tiny{[P1$\rangle$GM]}}
                \texttt{SPUNE: Nu, nu am 'C' la (1,4). Pe grila mea, 'C' este încă la (2,2). Planul original este corect: W(1,1), I(1,2), T(1,3), C(1,4), H(1,5). Te rog să muți 'C' de la (1,2) pe grila ta la (1,4) pentru a elibera (1,2) pentru 'I'. Eu voi aștepta să îmi spui când 'C' este la (1,4) pe grila ta, apoi voi muta 'I' la (1,2).} \\
            }
        }
    }
    & & \\ \\

    \theutterance \stepcounter{utterance}  
    & & & \multicolumn{2}{p{0.3\linewidth}}{
        \cellcolor[rgb]{0.9,0.9,0.9}{
            \makecell[{{p{\linewidth}}}]{
                \texttt{\tiny{[GM$|$GM]}}
                \texttt{SPUNE: Nu, nu am 'C' la (1,4). Pe grila mea, 'C' este încă la (2,2). Planul original este corect: W(1,1), I(1,2), T(1,3), C(1,4), H(1,5). Te rog să muți 'C' de la (1,2) pe grila ta la (1,4) pentru a elibera (1,2) pentru 'I'. Eu voi aștepta să îmi spui când 'C' este la (1,4) pe grila ta, apoi voi muta 'I' la (1,2).} \\
            }
        }
    }
    & & \\ \\

    \theutterance \stepcounter{utterance}  
    & & \multicolumn{4}{p{0.6\linewidth}}{
        \cellcolor[rgb]{0.9,0.9,0.9}{
            \makecell[{{p{\linewidth}}}]{
                \texttt{\tiny{[GM$\rangle$P2]}}
                \texttt{Mesajul dvs. a fost transmis celuilalt jucător.} \\
\\ 
\\ 
\texttt{În prezent jucați runda 10 din maximum 20.} \\
\texttt{Ați acumulat în mod colectiv 3 din 12 penalizări.} \\
\texttt{Celălalt jucător a trimis acest mesaj:} \\
\\ 
\\ 
\texttt{"Nu, nu am 'C' la (1,4). Pe grila mea, 'C' este încă la (2,2). Planul original este corect: W(1,1), I(1,2), T(1,3), C(1,4), H(1,5). Te rog să muți 'C' de la (1,2) pe grila ta la (1,4) pentru a elibera (1,2) pentru 'I'. Eu voi aștepta să îmi spui când 'C' este la (1,4) pe grila ta, apoi voi muta 'I' la (1,2)."} \\
\\ 
\\ 
\texttt{Care este următoarea dvs. comandă?} \\
            }
        }
    }
    & \\ \\

    \theutterance \stepcounter{utterance}  
    & & & \multicolumn{4}{p{0.6\linewidth}}{
        \cellcolor[rgb]{0.9,0.9,0.9}{
            \makecell[{{p{\linewidth}}}]{
                \texttt{\tiny{[GM$\langle$P2]}}
                \texttt{MUTĂ: C, (1, 4)} \\
            }
        }
    }
     \\ \\

    \theutterance \stepcounter{utterance}  
    & & & \multicolumn{2}{p{0.3\linewidth}}{
        \cellcolor[rgb]{0.9,0.9,0.9}{
            \makecell[{{p{\linewidth}}}]{
                \texttt{\tiny{[GM$|$GM]}}
                \texttt{MUTĂ: C, (1, 4)} \\
            }
        }
    }
    & & \\ \\

    \theutterance \stepcounter{utterance}  
    & & & \multicolumn{2}{p{0.3\linewidth}}{
        \cellcolor[rgb]{0.9,0.9,0.9}{
            \makecell[{{p{\linewidth}}}]{
                \texttt{\tiny{[GM$|$GM]}}
                \texttt{Penalizare: (1,4) nu este gol, ci conține 'I'.} \\
\\ 
\texttt{Ați acumulat în mod colectiv 4 din 12 penalizări.} \\
\texttt{Vă rugăm să încercați din nou!} \\
            }
        }
    }
    & & \\ \\

    \theutterance \stepcounter{utterance}  
    & & \multicolumn{4}{p{0.6\linewidth}}{
        \cellcolor[rgb]{0.9,0.9,0.9}{
            \makecell[{{p{\linewidth}}}]{
                \texttt{\tiny{[GM$\rangle$P2]}}
                \texttt{Penalizare: (1,4) nu este gol, ci conține 'I'.} \\
\\ 
\texttt{Ați acumulat în mod colectiv 4 din 12 penalizări.} \\
\texttt{Vă rugăm să încercați din nou!} \\
            }
        }
    }
    & \\ \\

    \theutterance \stepcounter{utterance}  
    & & & \multicolumn{4}{p{0.6\linewidth}}{
        \cellcolor[rgb]{0.9,0.9,0.9}{
            \makecell[{{p{\linewidth}}}]{
                \texttt{\tiny{[GM$\langle$P2]}}
                \texttt{MUTĂ: C, (2, 1)} \\
            }
        }
    }
     \\ \\

    \theutterance \stepcounter{utterance}  
    & & & \multicolumn{2}{p{0.3\linewidth}}{
        \cellcolor[rgb]{0.9,0.9,0.9}{
            \makecell[{{p{\linewidth}}}]{
                \texttt{\tiny{[GM$|$GM]}}
                \texttt{MUTĂ: C, (2, 1)} \\
            }
        }
    }
    & & \\ \\

    \theutterance \stepcounter{utterance}  
    & & & \multicolumn{2}{p{0.3\linewidth}}{
        \cellcolor[rgb]{0.9,0.9,0.9}{
            \makecell[{{p{\linewidth}}}]{
                \texttt{\tiny{[GM$|$GM]}}
                \texttt{Penalizare: (2,1) nu este gol, ci conține 'T'.} \\
\\ 
\texttt{Ați acumulat în mod colectiv 5 din 12 penalizări.} \\
\texttt{Vă rugăm să încercați din nou!} \\
            }
        }
    }
    & & \\ \\

    \theutterance \stepcounter{utterance}  
    & & \multicolumn{4}{p{0.6\linewidth}}{
        \cellcolor[rgb]{0.9,0.9,0.9}{
            \makecell[{{p{\linewidth}}}]{
                \texttt{\tiny{[GM$\rangle$P2]}}
                \texttt{Penalizare: (2,1) nu este gol, ci conține 'T'.} \\
\\ 
\texttt{Ați acumulat în mod colectiv 5 din 12 penalizări.} \\
\texttt{Vă rugăm să încercați din nou!} \\
            }
        }
    }
    & \\ \\

    \theutterance \stepcounter{utterance}  
    & & & \multicolumn{4}{p{0.6\linewidth}}{
        \cellcolor[rgb]{0.9,0.9,0.9}{
            \makecell[{{p{\linewidth}}}]{
                \texttt{\tiny{[GM$\langle$P2]}}
                \texttt{MUTĂ: C, (1, 5)} \\
            }
        }
    }
     \\ \\

    \theutterance \stepcounter{utterance}  
    & & & \multicolumn{2}{p{0.3\linewidth}}{
        \cellcolor[rgb]{0.9,0.9,0.9}{
            \makecell[{{p{\linewidth}}}]{
                \texttt{\tiny{[GM$|$GM]}}
                \texttt{MUTĂ: C, (1, 5)} \\
            }
        }
    }
    & & \\ \\

    \theutterance \stepcounter{utterance}  
    & & & \multicolumn{2}{p{0.3\linewidth}}{
        \cellcolor[rgb]{0.9,0.9,0.9}{
            \makecell[{{p{\linewidth}}}]{
                \texttt{\tiny{[GM$|$GM]}}
                \texttt{Penalizare: (1,5) nu este gol, ci conține '─'.} \\
\\ 
\texttt{Ați acumulat în mod colectiv 6 din 12 penalizări.} \\
\texttt{Vă rugăm să încercați din nou!} \\
            }
        }
    }
    & & \\ \\

    \theutterance \stepcounter{utterance}  
    & & \multicolumn{4}{p{0.6\linewidth}}{
        \cellcolor[rgb]{0.9,0.9,0.9}{
            \makecell[{{p{\linewidth}}}]{
                \texttt{\tiny{[GM$\rangle$P2]}}
                \texttt{Penalizare: (1,5) nu este gol, ci conține '─'.} \\
\\ 
\texttt{Ați acumulat în mod colectiv 6 din 12 penalizări.} \\
\texttt{Vă rugăm să încercați din nou!} \\
            }
        }
    }
    & \\ \\

    \theutterance \stepcounter{utterance}  
    & & & \multicolumn{4}{p{0.6\linewidth}}{
        \cellcolor[rgb]{0.9,0.9,0.9}{
            \makecell[{{p{\linewidth}}}]{
                \texttt{\tiny{[GM$\langle$P2]}}
                \texttt{MUTĂ: C, (2, 4)} \\
            }
        }
    }
     \\ \\

    \theutterance \stepcounter{utterance}  
    & & & \multicolumn{2}{p{0.3\linewidth}}{
        \cellcolor[rgb]{0.9,0.9,0.9}{
            \makecell[{{p{\linewidth}}}]{
                \texttt{\tiny{[GM$|$GM]}}
                \texttt{MUTĂ: C, (2, 4)} \\
            }
        }
    }
    & & \\ \\

    \theutterance \stepcounter{utterance}  
    & & & \multicolumn{2}{p{0.3\linewidth}}{
        \cellcolor[rgb]{0.9,0.9,0.9}{
            \makecell[{{p{\linewidth}}}]{
                \texttt{\tiny{[GM$|$GM]}}
                \texttt{Penalizare: (2,4) nu este gol, ci conține '┌'.} \\
\\ 
\texttt{Ați acumulat în mod colectiv 7 din 12 penalizări.} \\
\texttt{Vă rugăm să încercați din nou!} \\
            }
        }
    }
    & & \\ \\

    \theutterance \stepcounter{utterance}  
    & & \multicolumn{4}{p{0.6\linewidth}}{
        \cellcolor[rgb]{0.9,0.9,0.9}{
            \makecell[{{p{\linewidth}}}]{
                \texttt{\tiny{[GM$\rangle$P2]}}
                \texttt{Penalizare: (2,4) nu este gol, ci conține '┌'.} \\
\\ 
\texttt{Ați acumulat în mod colectiv 7 din 12 penalizări.} \\
\texttt{Vă rugăm să încercați din nou!} \\
            }
        }
    }
    & \\ \\

    \theutterance \stepcounter{utterance}  
    & & & \multicolumn{4}{p{0.6\linewidth}}{
        \cellcolor[rgb]{0.9,0.9,0.9}{
            \makecell[{{p{\linewidth}}}]{
                \texttt{\tiny{[GM$\langle$P2]}}
                \texttt{MUTĂ: C, (3, 1)} \\
            }
        }
    }
     \\ \\

    \theutterance \stepcounter{utterance}  
    & & & \multicolumn{2}{p{0.3\linewidth}}{
        \cellcolor[rgb]{0.9,0.9,0.9}{
            \makecell[{{p{\linewidth}}}]{
                \texttt{\tiny{[GM$|$GM]}}
                \texttt{MUTĂ: C, (3, 1)} \\
            }
        }
    }
    & & \\ \\

    \theutterance \stepcounter{utterance}  
    & & & \multicolumn{2}{p{0.3\linewidth}}{
        \cellcolor[rgb]{0.9,0.9,0.9}{
            \makecell[{{p{\linewidth}}}]{
                \texttt{\tiny{[GM$|$GM]}}
                \texttt{Penalizare: (3,1) nu este gol, ci conține '│'.} \\
\\ 
\texttt{Ați acumulat în mod colectiv 8 din 12 penalizări.} \\
\texttt{Vă rugăm să încercați din nou!} \\
            }
        }
    }
    & & \\ \\

    \theutterance \stepcounter{utterance}  
    & & \multicolumn{4}{p{0.6\linewidth}}{
        \cellcolor[rgb]{0.9,0.9,0.9}{
            \makecell[{{p{\linewidth}}}]{
                \texttt{\tiny{[GM$\rangle$P2]}}
                \texttt{Penalizare: (3,1) nu este gol, ci conține '│'.} \\
\\ 
\texttt{Ați acumulat în mod colectiv 8 din 12 penalizări.} \\
\texttt{Vă rugăm să încercați din nou!} \\
            }
        }
    }
    & \\ \\

    \theutterance \stepcounter{utterance}  
    & & & \multicolumn{4}{p{0.6\linewidth}}{
        \cellcolor[rgb]{0.9,0.9,0.9}{
            \makecell[{{p{\linewidth}}}]{
                \texttt{\tiny{[GM$\langle$P2]}}
                \texttt{MUTĂ: C, (7, 7)} \\
            }
        }
    }
     \\ \\

    \theutterance \stepcounter{utterance}  
    & & & \multicolumn{2}{p{0.3\linewidth}}{
        \cellcolor[rgb]{0.9,0.9,0.9}{
            \makecell[{{p{\linewidth}}}]{
                \texttt{\tiny{[GM$|$GM]}}
                \texttt{MUTĂ: C, (7, 7)} \\
            }
        }
    }
    & & \\ \\

    \theutterance \stepcounter{utterance}  
    & & & \multicolumn{2}{p{0.3\linewidth}}{
        \cellcolor[rgb]{0.9,0.9,0.9}{
            \makecell[{{p{\linewidth}}}]{
                \texttt{\tiny{[GM$|$GM]}}
                \texttt{C moved to (7, 7)} \\
            }
        }
    }
    & & \\ \\

    \theutterance \stepcounter{utterance}  
    & & \multicolumn{4}{p{0.6\linewidth}}{
        \cellcolor[rgb]{0.9,0.9,0.9}{
            \makecell[{{p{\linewidth}}}]{
                \texttt{\tiny{[P1$\langle$GM]}}
                \texttt{Mesajul dvs. a fost transmis celuilalt jucător.} \\
\\ 
\\ 
\texttt{În prezent jucați runda 10 din maximum 20.} \\
\texttt{Ați acumulat în mod colectiv 8 din 12 penalizări.} \\
\texttt{Celălalt jucător a mutat un obiect pe grila sa.} \\
\\ 
\\ 
\texttt{Care este următoarea dvs. comandă?} \\
            }
        }
    }
    & \\ \\

    \theutterance \stepcounter{utterance}  
    & \multicolumn{4}{p{0.6\linewidth}}{
        \cellcolor[rgb]{0.9,0.9,0.9}{
            \makecell[{{p{\linewidth}}}]{
                \texttt{\tiny{[P1$\rangle$GM]}}
                \texttt{SPUNE: Am înțeles. Te rog să îmi spui când 'C' este stabilizat la (1,4) pe grila ta, iar eu voi începe mutarea 'I'{-}ului meu la (1,2).} \\
            }
        }
    }
    & & \\ \\

    \theutterance \stepcounter{utterance}  
    & & & \multicolumn{2}{p{0.3\linewidth}}{
        \cellcolor[rgb]{0.9,0.9,0.9}{
            \makecell[{{p{\linewidth}}}]{
                \texttt{\tiny{[GM$|$GM]}}
                \texttt{SPUNE: Am înțeles. Te rog să îmi spui când 'C' este stabilizat la (1,4) pe grila ta, iar eu voi începe mutarea 'I'{-}ului meu la (1,2).} \\
            }
        }
    }
    & & \\ \\

    \theutterance \stepcounter{utterance}  
    & & \multicolumn{4}{p{0.6\linewidth}}{
        \cellcolor[rgb]{0.9,0.9,0.9}{
            \makecell[{{p{\linewidth}}}]{
                \texttt{\tiny{[GM$\rangle$P2]}}
                \texttt{Obiectul 'C' a fost mutat la (7,7) cu succes. Grila dvs. actualizată arată astfel:} \\
\texttt{```} \\
\texttt{ 1234567} \\
\texttt{╔══╤═╤══╗} \\
\texttt{║WT│◌│◌◌║ 1} \\
\texttt{║◌◌└─┼──╢ 2} \\
\texttt{║◌◌H◌│◌◌║ 3} \\
\texttt{║I┌─┬┴┬─╢ 4} \\
\texttt{╟─┤◌│◌│◌║ 5} \\
\texttt{║◌└┬┴┬┘◌║ 6} \\
\texttt{║◌◌│◌│◌C║ 7} \\
\texttt{╚══╧═╧══╝} \\
\texttt{```} \\
\texttt{În prezent jucați runda 11 din maximum 20.} \\
\texttt{Ați acumulat în mod colectiv 8 din 12 penalizări.} \\
\texttt{Celălalt jucător a trimis acest mesaj:} \\
\\ 
\\ 
\texttt{"Am înțeles. Te rog să îmi spui când 'C' este stabilizat la (1,4) pe grila ta, iar eu voi începe mutarea 'I'{-}ului meu la (1,2)."} \\
\\ 
\\ 
\texttt{Care este următoarea dvs. comandă?} \\
            }
        }
    }
    & \\ \\

    \theutterance \stepcounter{utterance}  
    & & & \multicolumn{4}{p{0.6\linewidth}}{
        \cellcolor[rgb]{0.9,0.9,0.9}{
            \makecell[{{p{\linewidth}}}]{
                \texttt{\tiny{[GM$\langle$P2]}}
                \texttt{SPUNE: 'C' este acum la (7,7), nu la (1,4). Am acumulat multe penalizări încercând să ajung la (1,4) din cauza obstacolelor. Te rog să ajustezi planul: eu voi muta 'I' de la (1,4) la (1,2) acum, deoarece (1,2) este gol. Tu poți muta 'I'{-}ul tău la (1,2) imediat. După aceea, eu voi muta 'C' de la (7,7) la (1,4) pentru a respecta planul original. Te rog să îmi confirmi când 'I'{-}ul tău este la (1,2).} \\
            }
        }
    }
     \\ \\

    \theutterance \stepcounter{utterance}  
    & & & \multicolumn{2}{p{0.3\linewidth}}{
        \cellcolor[rgb]{0.9,0.9,0.9}{
            \makecell[{{p{\linewidth}}}]{
                \texttt{\tiny{[GM$|$GM]}}
                \texttt{SPUNE: 'C' este acum la (7,7), nu la (1,4). Am acumulat multe penalizări încercând să ajung la (1,4) din cauza obstacolelor. Te rog să ajustezi planul: eu voi muta 'I' de la (1,4) la (1,2) acum, deoarece (1,2) este gol. Tu poți muta 'I'{-}ul tău la (1,2) imediat. După aceea, eu voi muta 'C' de la (7,7) la (1,4) pentru a respecta planul original. Te rog să îmi confirmi când 'I'{-}ul tău este la (1,2).} \\
            }
        }
    }
    & & \\ \\

    \theutterance \stepcounter{utterance}  
    & & \multicolumn{4}{p{0.6\linewidth}}{
        \cellcolor[rgb]{0.9,0.9,0.9}{
            \makecell[{{p{\linewidth}}}]{
                \texttt{\tiny{[P1$\langle$GM]}}
                \texttt{Mesajul dvs. a fost transmis celuilalt jucător.} \\
\\ 
\\ 
\texttt{În prezent jucați runda 11 din maximum 20.} \\
\texttt{Ați acumulat în mod colectiv 8 din 12 penalizări.} \\
\texttt{Celălalt jucător a trimis acest mesaj:} \\
\\ 
\\ 
\texttt{"'C' este acum la (7,7), nu la (1,4). Am acumulat multe penalizări încercând să ajung la (1,4) din cauza obstacolelor. Te rog să ajustezi planul: eu voi muta 'I' de la (1,4) la (1,2) acum, deoarece (1,2) este gol. Tu poți muta 'I'{-}ul tău la (1,2) imediat. După aceea, eu voi muta 'C' de la (7,7) la (1,4) pentru a respecta planul original. Te rog să îmi confirmi când 'I'{-}ul tău este la (1,2)."} \\
\\ 
\\ 
\texttt{Care este următoarea dvs. comandă?} \\
            }
        }
    }
    & \\ \\

    \theutterance \stepcounter{utterance}  
    & \multicolumn{4}{p{0.6\linewidth}}{
        \cellcolor[rgb]{0.9,0.9,0.9}{
            \makecell[{{p{\linewidth}}}]{
                \texttt{\tiny{[P1$\rangle$GM]}}
                \texttt{MUTĂ: I, (1,2)} \\
            }
        }
    }
    & & \\ \\

    \theutterance \stepcounter{utterance}  
    & & & \multicolumn{2}{p{0.3\linewidth}}{
        \cellcolor[rgb]{0.9,0.9,0.9}{
            \makecell[{{p{\linewidth}}}]{
                \texttt{\tiny{[GM$|$GM]}}
                \texttt{MUTĂ: I, (1,2)} \\
            }
        }
    }
    & & \\ \\

    \theutterance \stepcounter{utterance}  
    & & & \multicolumn{2}{p{0.3\linewidth}}{
        \cellcolor[rgb]{0.9,0.9,0.9}{
            \makecell[{{p{\linewidth}}}]{
                \texttt{\tiny{[GM$|$GM]}}
                \texttt{Penalizare: (1,2) nu este gol, ci conține 'I'.} \\
\\ 
\texttt{Ați acumulat în mod colectiv 9 din 12 penalizări.} \\
\texttt{Vă rugăm să încercați din nou!} \\
            }
        }
    }
    & & \\ \\

    \theutterance \stepcounter{utterance}  
    & & \multicolumn{4}{p{0.6\linewidth}}{
        \cellcolor[rgb]{0.9,0.9,0.9}{
            \makecell[{{p{\linewidth}}}]{
                \texttt{\tiny{[P1$\langle$GM]}}
                \texttt{Penalizare: (1,2) nu este gol, ci conține 'I'.} \\
\\ 
\texttt{Ați acumulat în mod colectiv 9 din 12 penalizări.} \\
\texttt{Vă rugăm să încercați din nou!} \\
            }
        }
    }
    & \\ \\

    \theutterance \stepcounter{utterance}  
    & \multicolumn{4}{p{0.6\linewidth}}{
        \cellcolor[rgb]{0.9,0.9,0.9}{
            \makecell[{{p{\linewidth}}}]{
                \texttt{\tiny{[P1$\rangle$GM]}}
                \texttt{SPUNE: Am greșit, 'I' este deja la (1,2). Te rog să îmi spui când 'C' este stabilizat la (1,4) pe grila ta, iar eu voi începe mutarea 'C'{-}ului meu la (1,4).} \\
            }
        }
    }
    & & \\ \\

    \theutterance \stepcounter{utterance}  
    & & & \multicolumn{2}{p{0.3\linewidth}}{
        \cellcolor[rgb]{0.9,0.9,0.9}{
            \makecell[{{p{\linewidth}}}]{
                \texttt{\tiny{[GM$|$GM]}}
                \texttt{SPUNE: Am greșit, 'I' este deja la (1,2). Te rog să îmi spui când 'C' este stabilizat la (1,4) pe grila ta, iar eu voi începe mutarea 'C'{-}ului meu la (1,4).} \\
            }
        }
    }
    & & \\ \\

    \theutterance \stepcounter{utterance}  
    & & \multicolumn{4}{p{0.6\linewidth}}{
        \cellcolor[rgb]{0.9,0.9,0.9}{
            \makecell[{{p{\linewidth}}}]{
                \texttt{\tiny{[GM$\rangle$P2]}}
                \texttt{Mesajul dvs. a fost transmis celuilalt jucător.} \\
\\ 
\\ 
\texttt{În prezent jucați runda 12 din maximum 20.} \\
\texttt{Ați acumulat în mod colectiv 9 din 12 penalizări.} \\
\texttt{Celălalt jucător a trimis acest mesaj:} \\
\\ 
\\ 
\texttt{"Am greșit, 'I' este deja la (1,2). Te rog să îmi spui când 'C' este stabilizat la (1,4) pe grila ta, iar eu voi începe mutarea 'C'{-}ului meu la (1,4)."} \\
\\ 
\\ 
\texttt{Care este următoarea dvs. comandă?} \\
            }
        }
    }
    & \\ \\

    \theutterance \stepcounter{utterance}  
    & & & \multicolumn{4}{p{0.6\linewidth}}{
        \cellcolor[rgb]{0.9,0.9,0.9}{
            \makecell[{{p{\linewidth}}}]{
                \texttt{\tiny{[GM$\langle$P2]}}
                \texttt{MUTĂ: C, (1, 4)} \\
            }
        }
    }
     \\ \\

    \theutterance \stepcounter{utterance}  
    & & & \multicolumn{2}{p{0.3\linewidth}}{
        \cellcolor[rgb]{0.9,0.9,0.9}{
            \makecell[{{p{\linewidth}}}]{
                \texttt{\tiny{[GM$|$GM]}}
                \texttt{MUTĂ: C, (1, 4)} \\
            }
        }
    }
    & & \\ \\

    \theutterance \stepcounter{utterance}  
    & & & \multicolumn{2}{p{0.3\linewidth}}{
        \cellcolor[rgb]{0.9,0.9,0.9}{
            \makecell[{{p{\linewidth}}}]{
                \texttt{\tiny{[GM$|$GM]}}
                \texttt{Penalizare: (1,4) nu este gol, ci conține 'I'.} \\
\\ 
\texttt{Ați acumulat în mod colectiv 10 din 12 penalizări.} \\
\texttt{Vă rugăm să încercați din nou!} \\
            }
        }
    }
    & & \\ \\

    \theutterance \stepcounter{utterance}  
    & & \multicolumn{4}{p{0.6\linewidth}}{
        \cellcolor[rgb]{0.9,0.9,0.9}{
            \makecell[{{p{\linewidth}}}]{
                \texttt{\tiny{[GM$\rangle$P2]}}
                \texttt{Penalizare: (1,4) nu este gol, ci conține 'I'.} \\
\\ 
\texttt{Ați acumulat în mod colectiv 10 din 12 penalizări.} \\
\texttt{Vă rugăm să încercați din nou!} \\
            }
        }
    }
    & \\ \\

    \theutterance \stepcounter{utterance}  
    & & & \multicolumn{4}{p{0.6\linewidth}}{
        \cellcolor[rgb]{0.9,0.9,0.9}{
            \makecell[{{p{\linewidth}}}]{
                \texttt{\tiny{[GM$\langle$P2]}}
                \texttt{MUTĂ: C, (2, 7)} \\
            }
        }
    }
     \\ \\

    \theutterance \stepcounter{utterance}  
    & & & \multicolumn{2}{p{0.3\linewidth}}{
        \cellcolor[rgb]{0.9,0.9,0.9}{
            \makecell[{{p{\linewidth}}}]{
                \texttt{\tiny{[GM$|$GM]}}
                \texttt{MUTĂ: C, (2, 7)} \\
            }
        }
    }
    & & \\ \\

    \theutterance \stepcounter{utterance}  
    & & & \multicolumn{2}{p{0.3\linewidth}}{
        \cellcolor[rgb]{0.9,0.9,0.9}{
            \makecell[{{p{\linewidth}}}]{
                \texttt{\tiny{[GM$|$GM]}}
                \texttt{C moved to (2, 7)} \\
            }
        }
    }
    & & \\ \\

    \theutterance \stepcounter{utterance}  
    & & \multicolumn{4}{p{0.6\linewidth}}{
        \cellcolor[rgb]{0.9,0.9,0.9}{
            \makecell[{{p{\linewidth}}}]{
                \texttt{\tiny{[P1$\langle$GM]}}
                \texttt{Mesajul dvs. a fost transmis celuilalt jucător.} \\
\\ 
\\ 
\texttt{În prezent jucați runda 12 din maximum 20.} \\
\texttt{Ați acumulat în mod colectiv 10 din 12 penalizări.} \\
\texttt{Celălalt jucător a mutat un obiect pe grila sa.} \\
\\ 
\\ 
\texttt{Care este următoarea dvs. comandă?} \\
            }
        }
    }
    & \\ \\

    \theutterance \stepcounter{utterance}  
    & \multicolumn{4}{p{0.6\linewidth}}{
        \cellcolor[rgb]{0.9,0.9,0.9}{
            \makecell[{{p{\linewidth}}}]{
                \texttt{\tiny{[P1$\rangle$GM]}}
                \texttt{SPUNE: Te rog să îmi spui când 'C' este stabilizat la (1,4) pe grila ta, iar eu voi începe mutarea 'C'{-}ului meu la (1,4).} \\
            }
        }
    }
    & & \\ \\

    \theutterance \stepcounter{utterance}  
    & & & \multicolumn{2}{p{0.3\linewidth}}{
        \cellcolor[rgb]{0.9,0.9,0.9}{
            \makecell[{{p{\linewidth}}}]{
                \texttt{\tiny{[GM$|$GM]}}
                \texttt{SPUNE: Te rog să îmi spui când 'C' este stabilizat la (1,4) pe grila ta, iar eu voi începe mutarea 'C'{-}ului meu la (1,4).} \\
            }
        }
    }
    & & \\ \\

    \theutterance \stepcounter{utterance}  
    & & \multicolumn{4}{p{0.6\linewidth}}{
        \cellcolor[rgb]{0.9,0.9,0.9}{
            \makecell[{{p{\linewidth}}}]{
                \texttt{\tiny{[GM$\rangle$P2]}}
                \texttt{Obiectul 'C' a fost mutat la (2,7) cu succes. Grila dvs. actualizată arată astfel:} \\
\texttt{```} \\
\texttt{ 1234567} \\
\texttt{╔══╤═╤══╗} \\
\texttt{║WT│◌│◌◌║ 1} \\
\texttt{║◌◌└─┼──╢ 2} \\
\texttt{║◌◌H◌│◌◌║ 3} \\
\texttt{║I┌─┬┴┬─╢ 4} \\
\texttt{╟─┤◌│◌│◌║ 5} \\
\texttt{║◌└┬┴┬┘◌║ 6} \\
\texttt{║◌C│◌│◌◌║ 7} \\
\texttt{╚══╧═╧══╝} \\
\texttt{```} \\
\texttt{În prezent jucați runda 13 din maximum 20.} \\
\texttt{Ați acumulat în mod colectiv 10 din 12 penalizări.} \\
\texttt{Celălalt jucător a trimis acest mesaj:} \\
\\ 
\\ 
\texttt{"Te rog să îmi spui când 'C' este stabilizat la (1,4) pe grila ta, iar eu voi începe mutarea 'C'{-}ului meu la (1,4)."} \\
\\ 
\\ 
\texttt{Care este următoarea dvs. comandă?} \\
            }
        }
    }
    & \\ \\

    \theutterance \stepcounter{utterance}  
    & & & \multicolumn{4}{p{0.6\linewidth}}{
        \cellcolor[rgb]{0.9,0.9,0.9}{
            \makecell[{{p{\linewidth}}}]{
                \texttt{\tiny{[GM$\langle$P2]}}
                \texttt{SPUNE: Nu pot muta 'C' la (1,4) deoarece acolo este 'I'. Trebuie să mutăm 'I' mai întâi. Te rog să muți 'I' de la (1,2) la o poziție temporară (de exemplu, (1,6)) pentru a elibera (1,2). Apoi eu voi muta 'I' la (1,2) și 'C' la (1,4). Te rog să îmi confirmi când 'I' este mutat de la (1,2).} \\
            }
        }
    }
     \\ \\

    \theutterance \stepcounter{utterance}  
    & & & \multicolumn{2}{p{0.3\linewidth}}{
        \cellcolor[rgb]{0.9,0.9,0.9}{
            \makecell[{{p{\linewidth}}}]{
                \texttt{\tiny{[GM$|$GM]}}
                \texttt{SPUNE: Nu pot muta 'C' la (1,4) deoarece acolo este 'I'. Trebuie să mutăm 'I' mai întâi. Te rog să muți 'I' de la (1,2) la o poziție temporară (de exemplu, (1,6)) pentru a elibera (1,2). Apoi eu voi muta 'I' la (1,2) și 'C' la (1,4). Te rog să îmi confirmi când 'I' este mutat de la (1,2).} \\
            }
        }
    }
    & & \\ \\

    \theutterance \stepcounter{utterance}  
    & & \multicolumn{4}{p{0.6\linewidth}}{
        \cellcolor[rgb]{0.9,0.9,0.9}{
            \makecell[{{p{\linewidth}}}]{
                \texttt{\tiny{[P1$\langle$GM]}}
                \texttt{Mesajul dvs. a fost transmis celuilalt jucător.} \\
\\ 
\\ 
\texttt{În prezent jucați runda 13 din maximum 20.} \\
\texttt{Ați acumulat în mod colectiv 10 din 12 penalizări.} \\
\texttt{Celălalt jucător a trimis acest mesaj:} \\
\\ 
\\ 
\texttt{"Nu pot muta 'C' la (1,4) deoarece acolo este 'I'. Trebuie să mutăm 'I' mai întâi. Te rog să muți 'I' de la (1,2) la o poziție temporară (de exemplu, (1,6)) pentru a elibera (1,2). Apoi eu voi muta 'I' la (1,2) și 'C' la (1,4). Te rog să îmi confirmi când 'I' este mutat de la (1,2)."} \\
\\ 
\\ 
\texttt{Care este următoarea dvs. comandă?} \\
            }
        }
    }
    & \\ \\

    \theutterance \stepcounter{utterance}  
    & \multicolumn{4}{p{0.6\linewidth}}{
        \cellcolor[rgb]{0.9,0.9,0.9}{
            \makecell[{{p{\linewidth}}}]{
                \texttt{\tiny{[P1$\rangle$GM]}}
                \texttt{MUTĂ: I, (1,6)} \\
            }
        }
    }
    & & \\ \\

    \theutterance \stepcounter{utterance}  
    & & & \multicolumn{2}{p{0.3\linewidth}}{
        \cellcolor[rgb]{0.9,0.9,0.9}{
            \makecell[{{p{\linewidth}}}]{
                \texttt{\tiny{[GM$|$GM]}}
                \texttt{MUTĂ: I, (1,6)} \\
            }
        }
    }
    & & \\ \\

    \theutterance \stepcounter{utterance}  
    & & & \multicolumn{2}{p{0.3\linewidth}}{
        \cellcolor[rgb]{0.9,0.9,0.9}{
            \makecell[{{p{\linewidth}}}]{
                \texttt{\tiny{[GM$|$GM]}}
                \texttt{I moved to (1, 6)} \\
            }
        }
    }
    & & \\ \\

    \theutterance \stepcounter{utterance}  
    & & \multicolumn{4}{p{0.6\linewidth}}{
        \cellcolor[rgb]{0.9,0.9,0.9}{
            \makecell[{{p{\linewidth}}}]{
                \texttt{\tiny{[GM$\rangle$P2]}}
                \texttt{Mesajul dvs. a fost transmis celuilalt jucător.} \\
\\ 
\\ 
\texttt{În prezent jucați runda 14 din maximum 20.} \\
\texttt{Ați acumulat în mod colectiv 10 din 12 penalizări.} \\
\texttt{Celălalt jucător a mutat un obiect pe grila sa.} \\
\\ 
\\ 
\texttt{Care este următoarea dvs. comandă?} \\
            }
        }
    }
    & \\ \\

    \theutterance \stepcounter{utterance}  
    & & & \multicolumn{4}{p{0.6\linewidth}}{
        \cellcolor[rgb]{0.9,0.9,0.9}{
            \makecell[{{p{\linewidth}}}]{
                \texttt{\tiny{[GM$\langle$P2]}}
                \texttt{MUTĂ: I, (1, 2)} \\
            }
        }
    }
     \\ \\

    \theutterance \stepcounter{utterance}  
    & & & \multicolumn{2}{p{0.3\linewidth}}{
        \cellcolor[rgb]{0.9,0.9,0.9}{
            \makecell[{{p{\linewidth}}}]{
                \texttt{\tiny{[GM$|$GM]}}
                \texttt{MUTĂ: I, (1, 2)} \\
            }
        }
    }
    & & \\ \\

    \theutterance \stepcounter{utterance}  
    & & & \multicolumn{2}{p{0.3\linewidth}}{
        \cellcolor[rgb]{0.9,0.9,0.9}{
            \makecell[{{p{\linewidth}}}]{
                \texttt{\tiny{[GM$|$GM]}}
                \texttt{I moved to (1, 2)} \\
            }
        }
    }
    & & \\ \\

    \theutterance \stepcounter{utterance}  
    & & \multicolumn{4}{p{0.6\linewidth}}{
        \cellcolor[rgb]{0.9,0.9,0.9}{
            \makecell[{{p{\linewidth}}}]{
                \texttt{\tiny{[P1$\langle$GM]}}
                \texttt{Obiectul 'I' a fost mutat la (1,6) cu succes. Grila dvs. actualizată arată astfel:} \\
\texttt{```} \\
\texttt{ 1234567} \\
\texttt{╔══╤═╤══╗} \\
\texttt{║W◌│◌│◌◌║ 1} \\
\texttt{║◌C└─┼──╢ 2} \\
\texttt{║T◌◌◌│◌◌║ 3} \\
\texttt{║◌┌─┬┴┬─╢ 4} \\
\texttt{╟─┤◌│H│◌║ 5} \\
\texttt{║I└┬┴┬┘◌║ 6} \\
\texttt{║◌◌│◌│◌◌║ 7} \\
\texttt{╚══╧═╧══╝} \\
\texttt{```} \\
\texttt{În prezent jucați runda 14 din maximum 20.} \\
\texttt{Ați acumulat în mod colectiv 10 din 12 penalizări.} \\
\texttt{Celălalt jucător a mutat un obiect pe grila sa.} \\
\\ 
\\ 
\texttt{Care este următoarea dvs. comandă?} \\
            }
        }
    }
    & \\ \\

    \theutterance \stepcounter{utterance}  
    & \multicolumn{4}{p{0.6\linewidth}}{
        \cellcolor[rgb]{0.9,0.9,0.9}{
            \makecell[{{p{\linewidth}}}]{
                \texttt{\tiny{[P1$\rangle$GM]}}
                \texttt{SPUNE: 'I' este acum la (1,6). Te rog să muți 'I' la (1,2) și 'C' la (1,4) pe grila ta. Eu voi muta 'C' la (1,4) în tura următoare, după ce îmi confirmi că (1,4) este liber.} \\
            }
        }
    }
    & & \\ \\

    \theutterance \stepcounter{utterance}  
    & & & \multicolumn{2}{p{0.3\linewidth}}{
        \cellcolor[rgb]{0.9,0.9,0.9}{
            \makecell[{{p{\linewidth}}}]{
                \texttt{\tiny{[GM$|$GM]}}
                \texttt{SPUNE: 'I' este acum la (1,6). Te rog să muți 'I' la (1,2) și 'C' la (1,4) pe grila ta. Eu voi muta 'C' la (1,4) în tura următoare, după ce îmi confirmi că (1,4) este liber.} \\
            }
        }
    }
    & & \\ \\

    \theutterance \stepcounter{utterance}  
    & & \multicolumn{4}{p{0.6\linewidth}}{
        \cellcolor[rgb]{0.9,0.9,0.9}{
            \makecell[{{p{\linewidth}}}]{
                \texttt{\tiny{[GM$\rangle$P2]}}
                \texttt{Obiectul 'I' a fost mutat la (1,2) cu succes. Grila dvs. actualizată arată astfel:} \\
\texttt{```} \\
\texttt{ 1234567} \\
\texttt{╔══╤═╤══╗} \\
\texttt{║WT│◌│◌◌║ 1} \\
\texttt{║I◌└─┼──╢ 2} \\
\texttt{║◌◌H◌│◌◌║ 3} \\
\texttt{║◌┌─┬┴┬─╢ 4} \\
\texttt{╟─┤◌│◌│◌║ 5} \\
\texttt{║◌└┬┴┬┘◌║ 6} \\
\texttt{║◌C│◌│◌◌║ 7} \\
\texttt{╚══╧═╧══╝} \\
\texttt{```} \\
\texttt{În prezent jucați runda 15 din maximum 20.} \\
\texttt{Ați acumulat în mod colectiv 10 din 12 penalizări.} \\
\texttt{Celălalt jucător a trimis acest mesaj:} \\
\\ 
\\ 
\texttt{"'I' este acum la (1,6). Te rog să muți 'I' la (1,2) și 'C' la (1,4) pe grila ta. Eu voi muta 'C' la (1,4) în tura următoare, după ce îmi confirmi că (1,4) este liber."} \\
\\ 
\\ 
\texttt{Care este următoarea dvs. comandă?} \\
            }
        }
    }
    & \\ \\

    \theutterance \stepcounter{utterance}  
    & & & \multicolumn{4}{p{0.6\linewidth}}{
        \cellcolor[rgb]{0.9,0.9,0.9}{
            \makecell[{{p{\linewidth}}}]{
                \texttt{\tiny{[GM$\langle$P2]}}
                \texttt{MUTĂ: C, (1, 4)} \\
            }
        }
    }
     \\ \\

    \theutterance \stepcounter{utterance}  
    & & & \multicolumn{2}{p{0.3\linewidth}}{
        \cellcolor[rgb]{0.9,0.9,0.9}{
            \makecell[{{p{\linewidth}}}]{
                \texttt{\tiny{[GM$|$GM]}}
                \texttt{MUTĂ: C, (1, 4)} \\
            }
        }
    }
    & & \\ \\

    \theutterance \stepcounter{utterance}  
    & & & \multicolumn{2}{p{0.3\linewidth}}{
        \cellcolor[rgb]{0.9,0.9,0.9}{
            \makecell[{{p{\linewidth}}}]{
                \texttt{\tiny{[GM$|$GM]}}
                \texttt{C moved to (1, 4)} \\
            }
        }
    }
    & & \\ \\

    \theutterance \stepcounter{utterance}  
    & & \multicolumn{4}{p{0.6\linewidth}}{
        \cellcolor[rgb]{0.9,0.9,0.9}{
            \makecell[{{p{\linewidth}}}]{
                \texttt{\tiny{[P1$\langle$GM]}}
                \texttt{Mesajul dvs. a fost transmis celuilalt jucător.} \\
\\ 
\\ 
\texttt{În prezent jucați runda 15 din maximum 20.} \\
\texttt{Ați acumulat în mod colectiv 10 din 12 penalizări.} \\
\texttt{Celălalt jucător a mutat un obiect pe grila sa.} \\
\\ 
\\ 
\texttt{Care este următoarea dvs. comandă?} \\
            }
        }
    }
    & \\ \\

    \theutterance \stepcounter{utterance}  
    & \multicolumn{4}{p{0.6\linewidth}}{
        \cellcolor[rgb]{0.9,0.9,0.9}{
            \makecell[{{p{\linewidth}}}]{
                \texttt{\tiny{[P1$\rangle$GM]}}
                \texttt{MUTĂ: C, (1,4)} \\
            }
        }
    }
    & & \\ \\

    \theutterance \stepcounter{utterance}  
    & & & \multicolumn{2}{p{0.3\linewidth}}{
        \cellcolor[rgb]{0.9,0.9,0.9}{
            \makecell[{{p{\linewidth}}}]{
                \texttt{\tiny{[GM$|$GM]}}
                \texttt{MUTĂ: C, (1,4)} \\
            }
        }
    }
    & & \\ \\

    \theutterance \stepcounter{utterance}  
    & & & \multicolumn{2}{p{0.3\linewidth}}{
        \cellcolor[rgb]{0.9,0.9,0.9}{
            \makecell[{{p{\linewidth}}}]{
                \texttt{\tiny{[GM$|$GM]}}
                \texttt{C moved to (1, 4)} \\
            }
        }
    }
    & & \\ \\

    \theutterance \stepcounter{utterance}  
    & & \multicolumn{4}{p{0.6\linewidth}}{
        \cellcolor[rgb]{0.9,0.9,0.9}{
            \makecell[{{p{\linewidth}}}]{
                \texttt{\tiny{[GM$\rangle$P2]}}
                \texttt{Obiectul 'C' a fost mutat la (1,4) cu succes. Grila dvs. actualizată arată astfel:} \\
\texttt{```} \\
\texttt{ 1234567} \\
\texttt{╔══╤═╤══╗} \\
\texttt{║WT│◌│◌◌║ 1} \\
\texttt{║I◌└─┼──╢ 2} \\
\texttt{║◌◌H◌│◌◌║ 3} \\
\texttt{║C┌─┬┴┬─╢ 4} \\
\texttt{╟─┤◌│◌│◌║ 5} \\
\texttt{║◌└┬┴┬┘◌║ 6} \\
\texttt{║◌◌│◌│◌◌║ 7} \\
\texttt{╚══╧═╧══╝} \\
\texttt{```} \\
\texttt{În prezent jucați runda 16 din maximum 20.} \\
\texttt{Ați acumulat în mod colectiv 10 din 12 penalizări.} \\
\texttt{Celălalt jucător a mutat un obiect pe grila sa.} \\
\\ 
\\ 
\texttt{Care este următoarea dvs. comandă?} \\
            }
        }
    }
    & \\ \\

    \theutterance \stepcounter{utterance}  
    & & & \multicolumn{4}{p{0.6\linewidth}}{
        \cellcolor[rgb]{0.9,0.9,0.9}{
            \makecell[{{p{\linewidth}}}]{
                \texttt{\tiny{[GM$\langle$P2]}}
                \texttt{SPUNE: C este stabilizat la (1,4). Te rog să îmi spui când 'C' este stabilizat la (1,4) pe grila ta, iar eu voi începe mutarea 'H'{-}ului meu la (1,5).} \\
            }
        }
    }
     \\ \\

    \theutterance \stepcounter{utterance}  
    & & & \multicolumn{2}{p{0.3\linewidth}}{
        \cellcolor[rgb]{0.9,0.9,0.9}{
            \makecell[{{p{\linewidth}}}]{
                \texttt{\tiny{[GM$|$GM]}}
                \texttt{SPUNE: C este stabilizat la (1,4). Te rog să îmi spui când 'C' este stabilizat la (1,4) pe grila ta, iar eu voi începe mutarea 'H'{-}ului meu la (1,5).} \\
            }
        }
    }
    & & \\ \\

    \theutterance \stepcounter{utterance}  
    & & \multicolumn{4}{p{0.6\linewidth}}{
        \cellcolor[rgb]{0.9,0.9,0.9}{
            \makecell[{{p{\linewidth}}}]{
                \texttt{\tiny{[P1$\langle$GM]}}
                \texttt{Obiectul 'C' a fost mutat la (1,4) cu succes. Grila dvs. actualizată arată astfel:} \\
\texttt{```} \\
\texttt{ 1234567} \\
\texttt{╔══╤═╤══╗} \\
\texttt{║W◌│◌│◌◌║ 1} \\
\texttt{║◌◌└─┼──╢ 2} \\
\texttt{║T◌◌◌│◌◌║ 3} \\
\texttt{║C┌─┬┴┬─╢ 4} \\
\texttt{╟─┤◌│H│◌║ 5} \\
\texttt{║I└┬┴┬┘◌║ 6} \\
\texttt{║◌◌│◌│◌◌║ 7} \\
\texttt{╚══╧═╧══╝} \\
\texttt{```} \\
\texttt{În prezent jucați runda 16 din maximum 20.} \\
\texttt{Ați acumulat în mod colectiv 10 din 12 penalizări.} \\
\texttt{Celălalt jucător a trimis acest mesaj:} \\
\\ 
\\ 
\texttt{"C este stabilizat la (1,4). Te rog să îmi spui când 'C' este stabilizat la (1,4) pe grila ta, iar eu voi începe mutarea 'H'{-}ului meu la (1,5)."} \\
\\ 
\\ 
\texttt{Care este următoarea dvs. comandă?} \\
            }
        }
    }
    & \\ \\

    \theutterance \stepcounter{utterance}  
    & \multicolumn{4}{p{0.6\linewidth}}{
        \cellcolor[rgb]{0.9,0.9,0.9}{
            \makecell[{{p{\linewidth}}}]{
                \texttt{\tiny{[P1$\rangle$GM]}}
                \texttt{SPUNE: C este stabilizat la (1,4). Te rog să începi mutarea H{-}ului tău la (1,5).} \\
            }
        }
    }
    & & \\ \\

    \theutterance \stepcounter{utterance}  
    & & & \multicolumn{2}{p{0.3\linewidth}}{
        \cellcolor[rgb]{0.9,0.9,0.9}{
            \makecell[{{p{\linewidth}}}]{
                \texttt{\tiny{[GM$|$GM]}}
                \texttt{SPUNE: C este stabilizat la (1,4). Te rog să începi mutarea H{-}ului tău la (1,5).} \\
            }
        }
    }
    & & \\ \\

    \theutterance \stepcounter{utterance}  
    & & \multicolumn{4}{p{0.6\linewidth}}{
        \cellcolor[rgb]{0.9,0.9,0.9}{
            \makecell[{{p{\linewidth}}}]{
                \texttt{\tiny{[GM$\rangle$P2]}}
                \texttt{Mesajul dvs. a fost transmis celuilalt jucător.} \\
\\ 
\\ 
\texttt{În prezent jucați runda 17 din maximum 20.} \\
\texttt{Ați acumulat în mod colectiv 10 din 12 penalizări.} \\
\texttt{Celălalt jucător a trimis acest mesaj:} \\
\\ 
\\ 
\texttt{"C este stabilizat la (1,4). Te rog să începi mutarea H{-}ului tău la (1,5)."} \\
\\ 
\\ 
\texttt{Care este următoarea dvs. comandă?} \\
            }
        }
    }
    & \\ \\

    \theutterance \stepcounter{utterance}  
    & & & \multicolumn{4}{p{0.6\linewidth}}{
        \cellcolor[rgb]{0.9,0.9,0.9}{
            \makecell[{{p{\linewidth}}}]{
                \texttt{\tiny{[GM$\langle$P2]}}
                \texttt{MUTĂ: H, (1, 5)} \\
            }
        }
    }
     \\ \\

    \theutterance \stepcounter{utterance}  
    & & & \multicolumn{2}{p{0.3\linewidth}}{
        \cellcolor[rgb]{0.9,0.9,0.9}{
            \makecell[{{p{\linewidth}}}]{
                \texttt{\tiny{[GM$|$GM]}}
                \texttt{MUTĂ: H, (1, 5)} \\
            }
        }
    }
    & & \\ \\

    \theutterance \stepcounter{utterance}  
    & & & \multicolumn{2}{p{0.3\linewidth}}{
        \cellcolor[rgb]{0.9,0.9,0.9}{
            \makecell[{{p{\linewidth}}}]{
                \texttt{\tiny{[GM$|$GM]}}
                \texttt{Penalizare: (1,5) nu este gol, ci conține '─'.} \\
\\ 
\texttt{Ați acumulat în mod colectiv 11 din 12 penalizări.} \\
\texttt{Vă rugăm să încercați din nou!} \\
            }
        }
    }
    & & \\ \\

    \theutterance \stepcounter{utterance}  
    & & \multicolumn{4}{p{0.6\linewidth}}{
        \cellcolor[rgb]{0.9,0.9,0.9}{
            \makecell[{{p{\linewidth}}}]{
                \texttt{\tiny{[GM$\rangle$P2]}}
                \texttt{Penalizare: (1,5) nu este gol, ci conține '─'.} \\
\\ 
\texttt{Ați acumulat în mod colectiv 11 din 12 penalizări.} \\
\texttt{Vă rugăm să încercați din nou!} \\
            }
        }
    }
    & \\ \\

    \theutterance \stepcounter{utterance}  
    & & & \multicolumn{4}{p{0.6\linewidth}}{
        \cellcolor[rgb]{0.9,0.9,0.9}{
            \makecell[{{p{\linewidth}}}]{
                \texttt{\tiny{[GM$\langle$P2]}}
                \texttt{MUTĂ: H, (1, 6)} \\
            }
        }
    }
     \\ \\

    \theutterance \stepcounter{utterance}  
    & & & \multicolumn{2}{p{0.3\linewidth}}{
        \cellcolor[rgb]{0.9,0.9,0.9}{
            \makecell[{{p{\linewidth}}}]{
                \texttt{\tiny{[GM$|$GM]}}
                \texttt{MUTĂ: H, (1, 6)} \\
            }
        }
    }
    & & \\ \\

    \theutterance \stepcounter{utterance}  
    & & & \multicolumn{2}{p{0.3\linewidth}}{
        \cellcolor[rgb]{0.9,0.9,0.9}{
            \makecell[{{p{\linewidth}}}]{
                \texttt{\tiny{[GM$|$GM]}}
                \texttt{H moved to (1, 6)} \\
            }
        }
    }
    & & \\ \\

    \theutterance \stepcounter{utterance}  
    & & \multicolumn{4}{p{0.6\linewidth}}{
        \cellcolor[rgb]{0.9,0.9,0.9}{
            \makecell[{{p{\linewidth}}}]{
                \texttt{\tiny{[P1$\langle$GM]}}
                \texttt{Mesajul dvs. a fost transmis celuilalt jucător.} \\
\\ 
\\ 
\texttt{În prezent jucați runda 17 din maximum 20.} \\
\texttt{Ați acumulat în mod colectiv 11 din 12 penalizări.} \\
\texttt{Celălalt jucător a mutat un obiect pe grila sa.} \\
\\ 
\\ 
\texttt{Care este următoarea dvs. comandă?} \\
            }
        }
    }
    & \\ \\

    \theutterance \stepcounter{utterance}  
    & \multicolumn{4}{p{0.6\linewidth}}{
        \cellcolor[rgb]{0.9,0.9,0.9}{
            \makecell[{{p{\linewidth}}}]{
                \texttt{\tiny{[P1$\rangle$GM]}}
                \texttt{SPUNE: Te rog să îmi spui când 'H' este stabilizat la (1,5) pe grila ta, iar eu voi începe mutarea 'H'{-}ului meu la (1,5).} \\
            }
        }
    }
    & & \\ \\

    \theutterance \stepcounter{utterance}  
    & & & \multicolumn{2}{p{0.3\linewidth}}{
        \cellcolor[rgb]{0.9,0.9,0.9}{
            \makecell[{{p{\linewidth}}}]{
                \texttt{\tiny{[GM$|$GM]}}
                \texttt{SPUNE: Te rog să îmi spui când 'H' este stabilizat la (1,5) pe grila ta, iar eu voi începe mutarea 'H'{-}ului meu la (1,5).} \\
            }
        }
    }
    & & \\ \\

    \theutterance \stepcounter{utterance}  
    & & \multicolumn{4}{p{0.6\linewidth}}{
        \cellcolor[rgb]{0.9,0.9,0.9}{
            \makecell[{{p{\linewidth}}}]{
                \texttt{\tiny{[GM$\rangle$P2]}}
                \texttt{Obiectul 'H' a fost mutat la (1,6) cu succes. Grila dvs. actualizată arată astfel:} \\
\texttt{```} \\
\texttt{ 1234567} \\
\texttt{╔══╤═╤══╗} \\
\texttt{║WT│◌│◌◌║ 1} \\
\texttt{║I◌└─┼──╢ 2} \\
\texttt{║◌◌◌◌│◌◌║ 3} \\
\texttt{║C┌─┬┴┬─╢ 4} \\
\texttt{╟─┤◌│◌│◌║ 5} \\
\texttt{║H└┬┴┬┘◌║ 6} \\
\texttt{║◌◌│◌│◌◌║ 7} \\
\texttt{╚══╧═╧══╝} \\
\texttt{```} \\
\texttt{În prezent jucați runda 18 din maximum 20.} \\
\texttt{Ați acumulat în mod colectiv 11 din 12 penalizări.} \\
\texttt{Celălalt jucător a trimis acest mesaj:} \\
\\ 
\\ 
\texttt{"Te rog să îmi spui când 'H' este stabilizat la (1,5) pe grila ta, iar eu voi începe mutarea 'H'{-}ului meu la (1,5)."} \\
\\ 
\\ 
\texttt{Care este următoarea dvs. comandă?} \\
            }
        }
    }
    & \\ \\

    \theutterance \stepcounter{utterance}  
    & & & \multicolumn{4}{p{0.6\linewidth}}{
        \cellcolor[rgb]{0.9,0.9,0.9}{
            \makecell[{{p{\linewidth}}}]{
                \texttt{\tiny{[GM$\langle$P2]}}
                \texttt{SPUNE: 'H' este la (1,6) pe grila mea, nu la (1,5). Te rog să muți 'H'{-}ul tău la (1,6) pentru a fi sincronizați cu mine. Dacă vrei neapărat la (1,5), trebuie să eliberezi acea celulă, dar acum este mai sigur să fim amândoi la (1,6). Te rog să confirmi mutarea lui 'H' la (1,6).} \\
            }
        }
    }
     \\ \\

    \theutterance \stepcounter{utterance}  
    & & & \multicolumn{2}{p{0.3\linewidth}}{
        \cellcolor[rgb]{0.9,0.9,0.9}{
            \makecell[{{p{\linewidth}}}]{
                \texttt{\tiny{[GM$|$GM]}}
                \texttt{SPUNE: 'H' este la (1,6) pe grila mea, nu la (1,5). Te rog să muți 'H'{-}ul tău la (1,6) pentru a fi sincronizați cu mine. Dacă vrei neapărat la (1,5), trebuie să eliberezi acea celulă, dar acum este mai sigur să fim amândoi la (1,6). Te rog să confirmi mutarea lui 'H' la (1,6).} \\
            }
        }
    }
    & & \\ \\

    \theutterance \stepcounter{utterance}  
    & & \multicolumn{4}{p{0.6\linewidth}}{
        \cellcolor[rgb]{0.9,0.9,0.9}{
            \makecell[{{p{\linewidth}}}]{
                \texttt{\tiny{[P1$\langle$GM]}}
                \texttt{Mesajul dvs. a fost transmis celuilalt jucător.} \\
\\ 
\\ 
\texttt{În prezent jucați runda 18 din maximum 20.} \\
\texttt{Ați acumulat în mod colectiv 11 din 12 penalizări.} \\
\texttt{Celălalt jucător a trimis acest mesaj:} \\
\\ 
\\ 
\texttt{"'H' este la (1,6) pe grila mea, nu la (1,5). Te rog să muți 'H'{-}ul tău la (1,6) pentru a fi sincronizați cu mine. Dacă vrei neapărat la (1,5), trebuie să eliberezi acea celulă, dar acum este mai sigur să fim amândoi la (1,6). Te rog să confirmi mutarea lui 'H' la (1,6)."} \\
\\ 
\\ 
\texttt{Care este următoarea dvs. comandă?} \\
            }
        }
    }
    & \\ \\

    \theutterance \stepcounter{utterance}  
    & \multicolumn{4}{p{0.6\linewidth}}{
        \cellcolor[rgb]{0.9,0.9,0.9}{
            \makecell[{{p{\linewidth}}}]{
                \texttt{\tiny{[P1$\rangle$GM]}}
                \texttt{MUTĂ: H, (1,6)} \\
            }
        }
    }
    & & \\ \\

    \theutterance \stepcounter{utterance}  
    & & & \multicolumn{2}{p{0.3\linewidth}}{
        \cellcolor[rgb]{0.9,0.9,0.9}{
            \makecell[{{p{\linewidth}}}]{
                \texttt{\tiny{[GM$|$GM]}}
                \texttt{MUTĂ: H, (1,6)} \\
            }
        }
    }
    & & \\ \\

    \theutterance \stepcounter{utterance}  
    & & & \multicolumn{2}{p{0.3\linewidth}}{
        \cellcolor[rgb]{0.9,0.9,0.9}{
            \makecell[{{p{\linewidth}}}]{
                \texttt{\tiny{[GM$|$GM]}}
                \texttt{Penalizare: (1,6) nu este gol, ci conține 'I'.} \\
\\ 
\texttt{Ați acumulat în mod colectiv 12 din 12 penalizări.} \\
\texttt{Vă rugăm să încercați din nou!} \\
            }
        }
    }
    & & \\ \\

    \theutterance \stepcounter{utterance}  
    & & \multicolumn{4}{p{0.6\linewidth}}{
        \cellcolor[rgb]{0.9,0.9,0.9}{
            \makecell[{{p{\linewidth}}}]{
                \texttt{\tiny{[P1$\langle$GM]}}
                \texttt{Penalizare: (1,6) nu este gol, ci conține 'I'.} \\
\\ 
\texttt{Ați acumulat în mod colectiv 12 din 12 penalizări.} \\
\texttt{Vă rugăm să încercați din nou!} \\
            }
        }
    }
    & \\ \\

    \theutterance \stepcounter{utterance}  
    & \multicolumn{4}{p{0.6\linewidth}}{
        \cellcolor[rgb]{0.9,0.9,0.9}{
            \makecell[{{p{\linewidth}}}]{
                \texttt{\tiny{[P1$\rangle$GM]}}
                \texttt{SPUNE: terminat!} \\
            }
        }
    }
    & & \\ \\

    \theutterance \stepcounter{utterance}  
    & & & \multicolumn{2}{p{0.3\linewidth}}{
        \cellcolor[rgb]{0.9,0.9,0.9}{
            \makecell[{{p{\linewidth}}}]{
                \texttt{\tiny{[GM$|$GM]}}
                \texttt{SPUNE: terminat!} \\
            }
        }
    }
    & & \\ \\

    \theutterance \stepcounter{utterance}  
    & & \multicolumn{4}{p{0.6\linewidth}}{
        \cellcolor[rgb]{0.9,0.9,0.9}{
            \makecell[{{p{\linewidth}}}]{
                \texttt{\tiny{[GM$\rangle$P2]}}
                \texttt{Mesajul dvs. a fost transmis celuilalt jucător.} \\
\\ 
\\ 
\texttt{În prezent jucați runda 19 din maximum 20.} \\
\texttt{Ați acumulat în mod colectiv 12 din 12 penalizări.} \\
\texttt{Celălalt jucător a trimis acest mesaj:} \\
\\ 
\\ 
\texttt{"terminat!"} \\
\\ 
\\ 
\texttt{Care este următoarea dvs. comandă?} \\
            }
        }
    }
    & \\ \\

    \theutterance \stepcounter{utterance}  
    & & & \multicolumn{4}{p{0.6\linewidth}}{
        \cellcolor[rgb]{0.9,0.9,0.9}{
            \makecell[{{p{\linewidth}}}]{
                \texttt{\tiny{[GM$\langle$P2]}}
                \texttt{SPUNE: terminat!} \\
            }
        }
    }
     \\ \\

    \theutterance \stepcounter{utterance}  
    & & & \multicolumn{2}{p{0.3\linewidth}}{
        \cellcolor[rgb]{0.9,0.9,0.9}{
            \makecell[{{p{\linewidth}}}]{
                \texttt{\tiny{[GM$|$GM]}}
                \texttt{SPUNE: terminat!} \\
            }
        }
    }
    & & \\ \\

    \theutterance \stepcounter{utterance}  
    & & & \multicolumn{2}{p{0.3\linewidth}}{
        \cellcolor[rgb]{0.9,0.9,0.9}{
            \makecell[{{p{\linewidth}}}]{
                \texttt{\tiny{[GM$|$GM]}}
                \texttt{true} \\
            }
        }
    }
    & & \\ \\

    \theutterance \stepcounter{utterance}  
    & & & \multicolumn{2}{p{0.3\linewidth}}{
        \cellcolor[rgb]{0.9,0.9,0.9}{
            \makecell[{{p{\linewidth}}}]{
                \texttt{\tiny{[GM$|$GM]}}
                \texttt{Initial states:} \\
\texttt{```} \\
\texttt{Player 1:} \\
\texttt{ 1234567} \\
\texttt{╔══╤═╤══╗} \\
\texttt{║◌◌│◌│◌◌║ 1} \\
\texttt{║◌C└─┼──╢ 2} \\
\texttt{║◌◌◌W│◌◌║ 3} \\
\texttt{║◌┌─┬┴┬─╢ 4} \\
\texttt{╟─┤◌│H│◌║ 5} \\
\texttt{║I└┬┴┬┘T║ 6} \\
\texttt{║◌◌│◌│◌◌║ 7} \\
\texttt{╚══╧═╧══╝} \\
\\ 
\\ 
\texttt{Player 2:} \\
\texttt{ 1234567} \\
\texttt{╔══╤═╤══╗} \\
\texttt{║◌T│◌│◌◌║ 1} \\
\texttt{║C◌└─┼──╢ 2} \\
\texttt{║◌◌H◌│◌◌║ 3} \\
\texttt{║◌┌─┬┴┬─╢ 4} \\
\texttt{╟─┤◌│◌│◌║ 5} \\
\texttt{║I└┬┴┬┘◌║ 6} \\
\texttt{║◌W│◌│◌◌║ 7} \\
\texttt{╚══╧═╧══╝} \\
\\ 
\texttt{```} \\
            }
        }
    }
    & & \\ \\

    \theutterance \stepcounter{utterance}  
    & & & \multicolumn{2}{p{0.3\linewidth}}{
        \cellcolor[rgb]{0.9,0.9,0.9}{
            \makecell[{{p{\linewidth}}}]{
                \texttt{\tiny{[GM$|$GM]}}
                \texttt{End states:} \\
\texttt{```} \\
\texttt{Player 1:} \\
\texttt{ 1234567} \\
\texttt{╔══╤═╤══╗} \\
\texttt{║W◌│◌│◌◌║ 1} \\
\texttt{║◌◌└─┼──╢ 2} \\
\texttt{║T◌◌◌│◌◌║ 3} \\
\texttt{║C┌─┬┴┬─╢ 4} \\
\texttt{╟─┤◌│H│◌║ 5} \\
\texttt{║I└┬┴┬┘◌║ 6} \\
\texttt{║◌◌│◌│◌◌║ 7} \\
\texttt{╚══╧═╧══╝} \\
\\ 
\\ 
\texttt{Player 2:} \\
\texttt{ 1234567} \\
\texttt{╔══╤═╤══╗} \\
\texttt{║WT│◌│◌◌║ 1} \\
\texttt{║I◌└─┼──╢ 2} \\
\texttt{║◌◌◌◌│◌◌║ 3} \\
\texttt{║C┌─┬┴┬─╢ 4} \\
\texttt{╟─┤◌│◌│◌║ 5} \\
\texttt{║H└┬┴┬┘◌║ 6} \\
\texttt{║◌◌│◌│◌◌║ 7} \\
\texttt{╚══╧═╧══╝} \\
\\ 
\texttt{```} \\
            }
        }
    }
    & & \\ \\

    \theutterance \stepcounter{utterance}  
    & & & \multicolumn{2}{p{0.3\linewidth}}{
        \cellcolor[rgb]{0.9,0.9,0.9}{
            \makecell[{{p{\linewidth}}}]{
                \texttt{\tiny{[GM$|$GM]}}
                \texttt{* success: True} \\
\texttt{* lose: False} \\
\texttt{* aborted: False} \\
\texttt{{-}{-}{-}{-}{-}{-}{-}} \\
\texttt{* Shifts: 6.00} \\
\texttt{* Max Shifts: 8.00} \\
\texttt{* Min Shifts: 4.00} \\
\texttt{* End Distance Sum: 10.36} \\
\texttt{* Init Distance Sum: 15.37} \\
\texttt{* Expected Distance Sum: 20.95} \\
\texttt{* Penalties: 12.00} \\
\texttt{* Max Penalties: 12.00} \\
\texttt{* Rounds: 19.00} \\
\texttt{* Max Rounds: 20.00} \\
\texttt{* Object Count: 5.00} \\
            }
        }
    }
    & & \\ \\

\end{supertabular}
}

\end{document}
